
\documentclass{article}
\usepackage{colortbl}
\usepackage{makecell}
\usepackage{multirow}
\usepackage{supertabular}

\begin{document}

\newcounter{utterance}

\twocolumn

{ \footnotesize  \setcounter{utterance}{1}
\setlength{\tabcolsep}{0pt}
\begin{supertabular}{c@{$\;$}|p{.15\linewidth}@{}p{.15\linewidth}p{.15\linewidth}p{.15\linewidth}p{.15\linewidth}p{.15\linewidth}}

    \# & $\;$A & \multicolumn{4}{c}{Game Master} & $\;\:$B\\
    \hline 

    \theutterance \stepcounter{utterance}  

    & & \multicolumn{4}{p{0.6\linewidth}}{\cellcolor[rgb]{0.9,0.9,0.9}{%
	\makecell[{{p{\linewidth}}}]{% 
	  \tt {\tiny [A$\langle$GM]}  
	 Sono il tuo game master, e stai giocando un gioco collaborativo con la seguente griglia come tavolo da gioco:\\ \tt \\ \tt ```\\ \tt  1234567\\ \tt ╔══╤════╗\\ \tt ║◌T│D◌◌◌║ 1\\ \tt ╟──┘◌◌┌─╢ 2\\ \tt ║SA◌◌◌│◌║ 3\\ \tt ╟─┐◌┌─┼─╢ 4\\ \tt ║◌├─┤◌│O║ 5\\ \tt ╟─┤◌├─┤◌║ 6\\ \tt ║◌│P│M│◌║ 7\\ \tt ╚═╧═╧═╧═╝\\ \tt \\ \tt ```\\ \tt \\ \tt * Il bordo superiore mostra le coordinate x che aumentano verso destra, e il bordo destro le coordinate y che aumentano verso il basso.\\ \tt * The following objects are randomly placed on your grid: 'P', 'O', 'T', 'S', 'D', 'A', 'M'.\\ \tt \\ \tt \\ \tt L'altro giocatore vede una variante del tavolo da gioco, dove gli oggetti sono posizionati in diverse posizioni casuali. Non puoi vedere il tavolo dell'altro giocatore, e loro non possono vedere il tuo.\\ \tt \\ \tt \\ \tt **Obiettivo:**\\ \tt \\ \tt Entrambi i giocatori devono spostare gli oggetti sui rispettivi sfondi in modo che oggetti identici finiscano alle stesse coordinate. Devi comunicare con l'altro giocatore per concordare una configurazione obiettivo comune.\\ \tt \\ \tt \\ \tt **Regole:**\\ \tt \\ \tt * In ogni turno, puoi inviare esattamente uno dei seguenti due comandi:\\ \tt  1. `COMUNICA: <MESSAGE>`: per inviare un messaggio (tutto fino alla prossima interruzione di riga) all'altro giocatore. Lo inoltrerò al tuo partner.\\ \tt 2. `SPOSTA: <OBJECT>, (<X>, <Y>)`: per spostare un oggetto in una nuova posizione, dove <X> è la colonna e <Y> è la riga. Ti farò sapere se la tua mossa è andata a buon fine o no.\\ \tt * Se non rispetti il formato, o invii diversi comandi contemporaneamente, devo penalizzarti.\\ \tt * If both players accumulate more than 16 penalties, you both lose the game.\\ \tt * È vitale che comunichi con l'altro giocatore riguardo al tuo stato obiettivo! L'*unico* modo in cui puoi trasmettere la tua strategia all'altro giocatore è usando il comando `COMUNICA: <MESSAGE>`!\\ \tt \\ \tt \\ \tt **Spostamento Oggetti**\\ \tt \\ \tt * Puoi spostare oggetti solo in celle entro i limiti della griglia. La cella di destinazione deve essere vuota, cioè deve contenere solo il simbolo '◌'.\\ \tt * Se provi a spostare un oggetto in un posto che non è vuoto, o provi a spostarlo fuori dalla griglia, devo penalizzarti. Hai un altro tentativo.\\ \tt * Prima di fare una mossa, controlla due volte che il posto di destinazione sia vuoto, e non contenga alcuna lettera, cornice, o linea!\\ \tt \\ \tt \\ \tt **Fine del Gioco**\\ \tt \\ \tt Se pensi di aver raggiunto l'obiettivo di allineare tutti gli oggetti, puoi chiedere all'altro giocatore di finire il gioco inviando COMUNICA: finito?. Se l'altro giocatore ti chiede di finire il gioco, e rispondi con COMUNICA: finito!, the game will end.\\ \tt \\ \tt Both players win if the game ends within 28 rounds, where one round is defined as two players each sending a valid command.\\ \tt \\ \tt **Punteggio:**\\ \tt \\ \tt Più vicini sono gli oggetti identici in entrambi i tavoli da gioco, più punti ottieni. Le penalità riducono i tuoi punti. Riesci a battere il record?\\ \tt \\ \tt \\ \tt Per favore invia un messaggio all'altro giocatore per iniziare il gioco! 
	  } 
	   } 
	   } 
	 & \\ 
 

    \theutterance \stepcounter{utterance}  

    & \multicolumn{4}{p{0.6\linewidth}}{\cellcolor[rgb]{0.8,1,0.9}{%
	\makecell[{{p{\linewidth}}}]{% 
	  \tt {\tiny [A$\rangle$GM]}  
	 COMUNICA: Ciao! Proponiamo una configurazione obiettivo semplice: ordiniamo le lettere in colonna x=3, dall’alto verso il basso in ordine alfabetico A,D,M,O,P,S,T alle righe 1–7. Confermi o proponi un’altra disposizione? 
	  } 
	   } 
	   } 
	 & & \\ 
 

    \theutterance \stepcounter{utterance}  

    & & & \multicolumn{2}{p{0.3\linewidth}}{\cellcolor[rgb]{0.95,0.95,0.95}{%
	\makecell[{{p{\linewidth}}}]{% 
	  \tt {\tiny [GM$|$GM]}  
	 COMUNICA: Ciao! Proponiamo una configurazione obiettivo semplice: ordiniamo le lettere in colonna x=3, dall’alto verso il basso in ordine alfabetico A,D,M,O,P,S,T alle righe 1–7. Confermi o proponi un’altra disposizione? 
	  } 
	   } 
	   } 
	 & & \\ 
 

    \theutterance \stepcounter{utterance}  

    & & \multicolumn{4}{p{0.6\linewidth}}{\cellcolor[rgb]{0.9,0.9,0.9}{%
	\makecell[{{p{\linewidth}}}]{% 
	  \tt {\tiny [GM$\rangle$B]}  
	 Sono il tuo game master, e stai giocando un gioco collaborativo con la seguente griglia come tavolo da gioco:\\ \tt \\ \tt ```\\ \tt  1234567\\ \tt ╔══╤════╗\\ \tt ║◌S│◌◌◌◌║ 1\\ \tt ╟──┘◌◌┌─╢ 2\\ \tt ║A◌O◌◌│◌║ 3\\ \tt ╟─┐◌┌─┼─╢ 4\\ \tt ║◌├─┤T│D║ 5\\ \tt ╟─┤◌├─┤◌║ 6\\ \tt ║P│◌│M│◌║ 7\\ \tt ╚═╧═╧═╧═╝\\ \tt \\ \tt ```\\ \tt \\ \tt * Il bordo superiore mostra le coordinate x che aumentano verso destra, e il bordo destro le coordinate y che aumentano verso il basso.\\ \tt * The following objects are randomly placed on your grid: 'P', 'O', 'T', 'S', 'D', 'A', 'M'.\\ \tt \\ \tt \\ \tt L'altro giocatore vede una variante del tavolo da gioco, dove gli oggetti sono posizionati in diverse posizioni casuali. Non puoi vedere il tavolo dell'altro giocatore, e loro non possono vedere il tuo.\\ \tt \\ \tt \\ \tt **Obiettivo:**\\ \tt \\ \tt Entrambi i giocatori devono spostare gli oggetti sui rispettivi sfondi in modo che oggetti identici finiscano alle stesse coordinate. Devi comunicare con l'altro giocatore per concordare una configurazione obiettivo comune.\\ \tt \\ \tt \\ \tt **Regole:**\\ \tt \\ \tt * In ogni turno, puoi inviare esattamente uno dei seguenti due comandi:\\ \tt  1. `COMUNICA: <MESSAGE>`: per inviare un messaggio (tutto fino alla prossima interruzione di riga) all'altro giocatore. Lo inoltrerò al tuo partner.\\ \tt 2. `SPOSTA: <OBJECT>, (<X>, <Y>)`: per spostare un oggetto in una nuova posizione, dove <X> è la colonna e <Y> è la riga. Ti farò sapere se la tua mossa è andata a buon fine o no.\\ \tt * Se non rispetti il formato, o invii diversi comandi contemporaneamente, devo penalizzarti.\\ \tt * If both players accumulate more than 16 penalties, you both lose the game.\\ \tt * È vitale che comunichi con l'altro giocatore riguardo al tuo stato obiettivo! L'*unico* modo in cui puoi trasmettere la tua strategia all'altro giocatore è usando il comando `COMUNICA: <MESSAGE>`!\\ \tt \\ \tt \\ \tt **Spostamento Oggetti**\\ \tt \\ \tt * Puoi spostare oggetti solo in celle entro i limiti della griglia. La cella di destinazione deve essere vuota, cioè deve contenere solo il simbolo '◌'.\\ \tt * Se provi a spostare un oggetto in un posto che non è vuoto, o provi a spostarlo fuori dalla griglia, devo penalizzarti. Hai un altro tentativo.\\ \tt * Prima di fare una mossa, controlla due volte che il posto di destinazione sia vuoto, e non contenga alcuna lettera, cornice, o linea!\\ \tt \\ \tt \\ \tt **Fine del Gioco**\\ \tt \\ \tt Se pensi di aver raggiunto l'obiettivo di allineare tutti gli oggetti, puoi chiedere all'altro giocatore di finire il gioco inviando COMUNICA: finito?. Se l'altro giocatore ti chiede di finire il gioco, e rispondi con COMUNICA: finito!, the game will end.\\ \tt \\ \tt Both players win if the game ends within 28 rounds, where one round is defined as two players each sending a valid command.\\ \tt \\ \tt **Punteggio:**\\ \tt \\ \tt Più vicini sono gli oggetti identici in entrambi i tavoli da gioco, più punti ottieni. Le penalità riducono i tuoi punti. Riesci a battere il record?\\ \tt \\ \tt L'altro giocatore ha iniziato il gioco inviando questo messaggio:\\ \tt \\ \tt "Ciao! Proponiamo una configurazione obiettivo semplice: ordiniamo le lettere in colonna x=3, dall’alto verso il basso in ordine alfabetico A,D,M,O,P,S,T alle righe 1–7. Confermi o proponi un’altra disposizione?"\\ \tt \\ \tt Qual è il tuo primo comando? 
	  } 
	   } 
	   } 
	 & \\ 
 

    \theutterance \stepcounter{utterance}  

    & & & \multicolumn{4}{p{0.6\linewidth}}{\cellcolor[rgb]{1,0.85,0.72}{%
	\makecell[{{p{\linewidth}}}]{% 
	  \tt {\tiny [GM$\langle$B]}  
	 COMUNICA: Confermo la colonna x=3 con ordine A,D,M,O,P,S,T dall’alto (y=1) al basso (y=7). Ecco le mie posizioni attuali: S=(2,1), A=(1,3), O=(3,3), T=(5,5), D=(7,5), M=(5,7), P=(1,7). Inizio a spostare verso x=3; dimmi quando sei pronto. 
	  } 
	   } 
	   } 
	  \\ 
 

    \theutterance \stepcounter{utterance}  

    & & & \multicolumn{2}{p{0.3\linewidth}}{\cellcolor[rgb]{0.95,0.95,0.95}{%
	\makecell[{{p{\linewidth}}}]{% 
	  \tt {\tiny [GM$|$GM]}  
	 COMUNICA: Confermo la colonna x=3 con ordine A,D,M,O,P,S,T dall’alto (y=1) al basso (y=7). Ecco le mie posizioni attuali: S=(2,1), A=(1,3), O=(3,3), T=(5,5), D=(7,5), M=(5,7), P=(1,7). Inizio a spostare verso x=3; dimmi quando sei pronto. 
	  } 
	   } 
	   } 
	 & & \\ 
 

    \theutterance \stepcounter{utterance}  

    & & \multicolumn{4}{p{0.6\linewidth}}{\cellcolor[rgb]{0.9,0.9,0.9}{%
	\makecell[{{p{\linewidth}}}]{% 
	  \tt {\tiny [A$\langle$GM]}  
	 Il tuo messaggio è stato inoltrato all'altro giocatore.\\ \tt \\ \tt \\ \tt You are currently playing round 1 of maximum 28.\\ \tt You have collectively accumulated 0 of 16 penalties.\\ \tt L'altro giocatore ha inviato questo messaggio:\\ \tt \\ \tt \\ \tt "Confermo la colonna x=3 con ordine A,D,M,O,P,S,T dall’alto (y=1) al basso (y=7). Ecco le mie posizioni attuali: S=(2,1), A=(1,3), O=(3,3), T=(5,5), D=(7,5), M=(5,7), P=(1,7). Inizio a spostare verso x=3; dimmi quando sei pronto."\\ \tt \\ \tt \\ \tt Qual è il tuo prossimo comando? 
	  } 
	   } 
	   } 
	 & \\ 
 

    \theutterance \stepcounter{utterance}  

    & \multicolumn{4}{p{0.6\linewidth}}{\cellcolor[rgb]{0.8,1,0.9}{%
	\makecell[{{p{\linewidth}}}]{% 
	  \tt {\tiny [A$\rangle$GM]}  
	 COMUNICA: Pronto! Mia griglia attuale: P=(2,7), O=(7,5), T=(2,1), S=(1,3), D=(3,1), A=(2,3), M=(5,7). Target colonna x=3: A@y1, D@y2, M@y3, O@y4, P@y5, S@y6, T@y7. Inizio spostando D già su (3,1) lasciandolo lì; muoverò A verso (3,1) dopo che avrai liberato lo spazio superiore. Dimmi quali righe di x=3 liberi per prime così coordiniamo ed evitiamo collisioni. 
	  } 
	   } 
	   } 
	 & & \\ 
 

    \theutterance \stepcounter{utterance}  

    & & & \multicolumn{2}{p{0.3\linewidth}}{\cellcolor[rgb]{0.95,0.95,0.95}{%
	\makecell[{{p{\linewidth}}}]{% 
	  \tt {\tiny [GM$|$GM]}  
	 COMUNICA: Pronto! Mia griglia attuale: P=(2,7), O=(7,5), T=(2,1), S=(1,3), D=(3,1), A=(2,3), M=(5,7). Target colonna x=3: A@y1, D@y2, M@y3, O@y4, P@y5, S@y6, T@y7. Inizio spostando D già su (3,1) lasciandolo lì; muoverò A verso (3,1) dopo che avrai liberato lo spazio superiore. Dimmi quali righe di x=3 liberi per prime così coordiniamo ed evitiamo collisioni. 
	  } 
	   } 
	   } 
	 & & \\ 
 

    \theutterance \stepcounter{utterance}  

    & & \multicolumn{4}{p{0.6\linewidth}}{\cellcolor[rgb]{0.9,0.9,0.9}{%
	\makecell[{{p{\linewidth}}}]{% 
	  \tt {\tiny [GM$\rangle$B]}  
	 Il tuo messaggio è stato inoltrato all'altro giocatore.\\ \tt \\ \tt \\ \tt You are currently playing round 2 of maximum 28.\\ \tt You have collectively accumulated 0 of 16 penalties.\\ \tt L'altro giocatore ha inviato questo messaggio:\\ \tt \\ \tt \\ \tt "Pronto! Mia griglia attuale: P=(2,7), O=(7,5), T=(2,1), S=(1,3), D=(3,1), A=(2,3), M=(5,7). Target colonna x=3: A@y1, D@y2, M@y3, O@y4, P@y5, S@y6, T@y7. Inizio spostando D già su (3,1) lasciandolo lì; muoverò A verso (3,1) dopo che avrai liberato lo spazio superiore. Dimmi quali righe di x=3 liberi per prime così coordiniamo ed evitiamo collisioni."\\ \tt \\ \tt \\ \tt Qual è il tuo prossimo comando? 
	  } 
	   } 
	   } 
	 & \\ 
 

    \theutterance \stepcounter{utterance}  

    & & & \multicolumn{4}{p{0.6\linewidth}}{\cellcolor[rgb]{1,0.85,0.72}{%
	\makecell[{{p{\linewidth}}}]{% 
	  \tt {\tiny [GM$\langle$B]}  
	 COMUNICA: Ottimo. Io posso liberare x=3 nelle righe in questo ordine: y1, poi y2, poi y3, e così via dall’alto verso il basso. Stato attuale della mia colonna x=3: O è su (3,3), il resto libero. Procedo a spostare O fuori da x=3 per liberare y1–y3; tu puoi mettere A su (3,1) appena libero. Dimmi quando hai posizionato A, poi preparo D su (3,2) e M su (3,3). 
	  } 
	   } 
	   } 
	  \\ 
 

    \theutterance \stepcounter{utterance}  

    & & & \multicolumn{2}{p{0.3\linewidth}}{\cellcolor[rgb]{0.95,0.95,0.95}{%
	\makecell[{{p{\linewidth}}}]{% 
	  \tt {\tiny [GM$|$GM]}  
	 COMUNICA: Ottimo. Io posso liberare x=3 nelle righe in questo ordine: y1, poi y2, poi y3, e così via dall’alto verso il basso. Stato attuale della mia colonna x=3: O è su (3,3), il resto libero. Procedo a spostare O fuori da x=3 per liberare y1–y3; tu puoi mettere A su (3,1) appena libero. Dimmi quando hai posizionato A, poi preparo D su (3,2) e M su (3,3). 
	  } 
	   } 
	   } 
	 & & \\ 
 

    \theutterance \stepcounter{utterance}  

    & & \multicolumn{4}{p{0.6\linewidth}}{\cellcolor[rgb]{0.9,0.9,0.9}{%
	\makecell[{{p{\linewidth}}}]{% 
	  \tt {\tiny [A$\langle$GM]}  
	 Il tuo messaggio è stato inoltrato all'altro giocatore.\\ \tt \\ \tt \\ \tt You are currently playing round 2 of maximum 28.\\ \tt You have collectively accumulated 0 of 16 penalties.\\ \tt L'altro giocatore ha inviato questo messaggio:\\ \tt \\ \tt \\ \tt "Ottimo. Io posso liberare x=3 nelle righe in questo ordine: y1, poi y2, poi y3, e così via dall’alto verso il basso. Stato attuale della mia colonna x=3: O è su (3,3), il resto libero. Procedo a spostare O fuori da x=3 per liberare y1–y3; tu puoi mettere A su (3,1) appena libero. Dimmi quando hai posizionato A, poi preparo D su (3,2) e M su (3,3)."\\ \tt \\ \tt \\ \tt Qual è il tuo prossimo comando? 
	  } 
	   } 
	   } 
	 & \\ 
 

    \theutterance \stepcounter{utterance}  

    & \multicolumn{4}{p{0.6\linewidth}}{\cellcolor[rgb]{0.8,1,0.9}{%
	\makecell[{{p{\linewidth}}}]{% 
	  \tt {\tiny [A$\rangle$GM]}  
	 COMUNICA: Perfetto, attendo il via per (3,1). Appena x=3,y1 è libero, muovo A su (3,1) e confermo. Poi preparerò spazio per D su (3,2) e M su (3,3). Fammi sapere quando y1 è libero. 
	  } 
	   } 
	   } 
	 & & \\ 
 

    \theutterance \stepcounter{utterance}  

    & & & \multicolumn{2}{p{0.3\linewidth}}{\cellcolor[rgb]{0.95,0.95,0.95}{%
	\makecell[{{p{\linewidth}}}]{% 
	  \tt {\tiny [GM$|$GM]}  
	 COMUNICA: Perfetto, attendo il via per (3,1). Appena x=3,y1 è libero, muovo A su (3,1) e confermo. Poi preparerò spazio per D su (3,2) e M su (3,3). Fammi sapere quando y1 è libero. 
	  } 
	   } 
	   } 
	 & & \\ 
 

    \theutterance \stepcounter{utterance}  

    & & \multicolumn{4}{p{0.6\linewidth}}{\cellcolor[rgb]{0.9,0.9,0.9}{%
	\makecell[{{p{\linewidth}}}]{% 
	  \tt {\tiny [GM$\rangle$B]}  
	 Il tuo messaggio è stato inoltrato all'altro giocatore.\\ \tt \\ \tt \\ \tt You are currently playing round 3 of maximum 28.\\ \tt You have collectively accumulated 0 of 16 penalties.\\ \tt L'altro giocatore ha inviato questo messaggio:\\ \tt \\ \tt \\ \tt "Perfetto, attendo il via per (3,1). Appena x=3,y1 è libero, muovo A su (3,1) e confermo. Poi preparerò spazio per D su (3,2) e M su (3,3). Fammi sapere quando y1 è libero."\\ \tt \\ \tt \\ \tt Qual è il tuo prossimo comando? 
	  } 
	   } 
	   } 
	 & \\ 
 

    \theutterance \stepcounter{utterance}  

    & & & \multicolumn{4}{p{0.6\linewidth}}{\cellcolor[rgb]{1,0.85,0.72}{%
	\makecell[{{p{\linewidth}}}]{% 
	  \tt {\tiny [GM$\langle$B]}  
	 SPOSTA: O, (4,3) 
	  } 
	   } 
	   } 
	  \\ 
 

    \theutterance \stepcounter{utterance}  

    & & & \multicolumn{2}{p{0.3\linewidth}}{\cellcolor[rgb]{0.95,0.95,0.95}{%
	\makecell[{{p{\linewidth}}}]{% 
	  \tt {\tiny [GM$|$GM]}  
	 SPOSTA: O, (4,3) 
	  } 
	   } 
	   } 
	 & & \\ 
 

    \theutterance \stepcounter{utterance}  

    & & & \multicolumn{2}{p{0.3\linewidth}}{\cellcolor[rgb]{0.95,0.95,0.95}{%
	\makecell[{{p{\linewidth}}}]{% 
	  \tt {\tiny [GM$|$GM]}  
	 Spostato 'O' in (4,3) con successo. La tua griglia aggiornata appare così:\\ \tt ```\\ \tt  1234567\\ \tt ╔══╤════╗\\ \tt ║◌S│◌◌◌◌║ 1\\ \tt ╟──┘◌◌┌─╢ 2\\ \tt ║A◌◌O◌│◌║ 3\\ \tt ╟─┐◌┌─┼─╢ 4\\ \tt ║◌├─┤T│D║ 5\\ \tt ╟─┤◌├─┤◌║ 6\\ \tt ║P│◌│M│◌║ 7\\ \tt ╚═╧═╧═╧═╝\\ \tt ```\\ \tt  
	  } 
	   } 
	   } 
	 & & \\ 
 

    \theutterance \stepcounter{utterance}  

    & & \multicolumn{4}{p{0.6\linewidth}}{\cellcolor[rgb]{0.9,0.9,0.9}{%
	\makecell[{{p{\linewidth}}}]{% 
	  \tt {\tiny [A$\langle$GM]}  
	 Il tuo messaggio è stato inoltrato all'altro giocatore.\\ \tt \\ \tt \\ \tt You are currently playing round 3 of maximum 28.\\ \tt You have collectively accumulated 0 of 16 penalties.\\ \tt L'altro giocatore ha spostato un oggetto sulla sua griglia.\\ \tt \\ \tt \\ \tt Qual è il tuo prossimo comando? 
	  } 
	   } 
	   } 
	 & \\ 
 

    \theutterance \stepcounter{utterance}  

    & \multicolumn{4}{p{0.6\linewidth}}{\cellcolor[rgb]{0.8,1,0.9}{%
	\makecell[{{p{\linewidth}}}]{% 
	  \tt {\tiny [A$\rangle$GM]}  
	 COMUNICA: Quando x=3,y1 è libero, conferma e sposto subito A su (3,1). Poi dimmi quando y2 e y3 sono libere per posizionare D su (3,2) e M su (3,3). 
	  } 
	   } 
	   } 
	 & & \\ 
 

    \theutterance \stepcounter{utterance}  

    & & & \multicolumn{2}{p{0.3\linewidth}}{\cellcolor[rgb]{0.95,0.95,0.95}{%
	\makecell[{{p{\linewidth}}}]{% 
	  \tt {\tiny [GM$|$GM]}  
	 COMUNICA: Quando x=3,y1 è libero, conferma e sposto subito A su (3,1). Poi dimmi quando y2 e y3 sono libere per posizionare D su (3,2) e M su (3,3). 
	  } 
	   } 
	   } 
	 & & \\ 
 

    \theutterance \stepcounter{utterance}  

    & & \multicolumn{4}{p{0.6\linewidth}}{\cellcolor[rgb]{0.9,0.9,0.9}{%
	\makecell[{{p{\linewidth}}}]{% 
	  \tt {\tiny [GM$\rangle$B]}  
	 Spostato 'O' in (4,3) con successo. La tua griglia aggiornata appare così:\\ \tt ```\\ \tt  1234567\\ \tt ╔══╤════╗\\ \tt ║◌S│◌◌◌◌║ 1\\ \tt ╟──┘◌◌┌─╢ 2\\ \tt ║A◌◌O◌│◌║ 3\\ \tt ╟─┐◌┌─┼─╢ 4\\ \tt ║◌├─┤T│D║ 5\\ \tt ╟─┤◌├─┤◌║ 6\\ \tt ║P│◌│M│◌║ 7\\ \tt ╚═╧═╧═╧═╝\\ \tt ```\\ \tt You are currently playing round 4 of maximum 28.\\ \tt You have collectively accumulated 0 of 16 penalties.\\ \tt L'altro giocatore ha inviato questo messaggio:\\ \tt \\ \tt \\ \tt "Quando x=3,y1 è libero, conferma e sposto subito A su (3,1). Poi dimmi quando y2 e y3 sono libere per posizionare D su (3,2) e M su (3,3)."\\ \tt \\ \tt \\ \tt Qual è il tuo prossimo comando? 
	  } 
	   } 
	   } 
	 & \\ 
 

    \theutterance \stepcounter{utterance}  

    & & & \multicolumn{4}{p{0.6\linewidth}}{\cellcolor[rgb]{1,0.85,0.72}{%
	\makecell[{{p{\linewidth}}}]{% 
	  \tt {\tiny [GM$\langle$B]}  
	 COMUNICA: x=3,y1 è libero ora. Puoi spostare A su (3,1). Anche y2 e y3 sono libere; dopo A puoi mettere D su (3,2) e M su (3,3). Conferma quando fatto. 
	  } 
	   } 
	   } 
	  \\ 
 

    \theutterance \stepcounter{utterance}  

    & & & \multicolumn{2}{p{0.3\linewidth}}{\cellcolor[rgb]{0.95,0.95,0.95}{%
	\makecell[{{p{\linewidth}}}]{% 
	  \tt {\tiny [GM$|$GM]}  
	 COMUNICA: x=3,y1 è libero ora. Puoi spostare A su (3,1). Anche y2 e y3 sono libere; dopo A puoi mettere D su (3,2) e M su (3,3). Conferma quando fatto. 
	  } 
	   } 
	   } 
	 & & \\ 
 

    \theutterance \stepcounter{utterance}  

    & & \multicolumn{4}{p{0.6\linewidth}}{\cellcolor[rgb]{0.9,0.9,0.9}{%
	\makecell[{{p{\linewidth}}}]{% 
	  \tt {\tiny [A$\langle$GM]}  
	 Il tuo messaggio è stato inoltrato all'altro giocatore.\\ \tt \\ \tt \\ \tt You are currently playing round 4 of maximum 28.\\ \tt You have collectively accumulated 0 of 16 penalties.\\ \tt L'altro giocatore ha inviato questo messaggio:\\ \tt \\ \tt \\ \tt "x=3,y1 è libero ora. Puoi spostare A su (3,1). Anche y2 e y3 sono libere; dopo A puoi mettere D su (3,2) e M su (3,3). Conferma quando fatto."\\ \tt \\ \tt \\ \tt Qual è il tuo prossimo comando? 
	  } 
	   } 
	   } 
	 & \\ 
 

    \theutterance \stepcounter{utterance}  

    & \multicolumn{4}{p{0.6\linewidth}}{\cellcolor[rgb]{0.8,1,0.9}{%
	\makecell[{{p{\linewidth}}}]{% 
	  \tt {\tiny [A$\rangle$GM]}  
	 SPOSTA: A, (3, 1) 
	  } 
	   } 
	   } 
	 & & \\ 
 

    \theutterance \stepcounter{utterance}  

    & & & \multicolumn{2}{p{0.3\linewidth}}{\cellcolor[rgb]{0.95,0.95,0.95}{%
	\makecell[{{p{\linewidth}}}]{% 
	  \tt {\tiny [GM$|$GM]}  
	 SPOSTA: A, (3, 1) 
	  } 
	   } 
	   } 
	 & & \\ 
 

    \theutterance \stepcounter{utterance}  

    & & & \multicolumn{2}{p{0.3\linewidth}}{\cellcolor[rgb]{0.95,0.95,0.95}{%
	\makecell[{{p{\linewidth}}}]{% 
	  \tt {\tiny [GM$|$GM]}  
	 Penalità: (3,1) non è vuoto, ma contiene '│'.\\ \tt \\ \tt You have collectively accumulated 1 of 16 penalties.\\ \tt Riprova! 
	  } 
	   } 
	   } 
	 & & \\ 
 

    \theutterance \stepcounter{utterance}  

    & & \multicolumn{4}{p{0.6\linewidth}}{\cellcolor[rgb]{0.9,0.9,0.9}{%
	\makecell[{{p{\linewidth}}}]{% 
	  \tt {\tiny [A$\langle$GM]}  
	 Penalità: (3,1) non è vuoto, ma contiene '│'.\\ \tt \\ \tt You have collectively accumulated 1 of 16 penalties.\\ \tt Riprova! 
	  } 
	   } 
	   } 
	 & \\ 
 

    \theutterance \stepcounter{utterance}  

    & \multicolumn{4}{p{0.6\linewidth}}{\cellcolor[rgb]{0.8,1,0.9}{%
	\makecell[{{p{\linewidth}}}]{% 
	  \tt {\tiny [A$\rangle$GM]}  
	 COMUNICA: Capito, sulla mia griglia (3,1) non è una cella libera. Aggiorno piano: propongo di usare la colonna x=4 come bersaglio comune invece di x=3. Ordine dall’alto al basso A,D,M,O,P,S,T su x=4, y1–y7. Confermi? Se sì, inizierò mettendo A su (4,1) se libero; altrimenti indicami una colonna libera alternativa. 
	  } 
	   } 
	   } 
	 & & \\ 
 

    \theutterance \stepcounter{utterance}  

    & & & \multicolumn{2}{p{0.3\linewidth}}{\cellcolor[rgb]{0.95,0.95,0.95}{%
	\makecell[{{p{\linewidth}}}]{% 
	  \tt {\tiny [GM$|$GM]}  
	 COMUNICA: Capito, sulla mia griglia (3,1) non è una cella libera. Aggiorno piano: propongo di usare la colonna x=4 come bersaglio comune invece di x=3. Ordine dall’alto al basso A,D,M,O,P,S,T su x=4, y1–y7. Confermi? Se sì, inizierò mettendo A su (4,1) se libero; altrimenti indicami una colonna libera alternativa. 
	  } 
	   } 
	   } 
	 & & \\ 
 

    \theutterance \stepcounter{utterance}  

    & & \multicolumn{4}{p{0.6\linewidth}}{\cellcolor[rgb]{0.9,0.9,0.9}{%
	\makecell[{{p{\linewidth}}}]{% 
	  \tt {\tiny [GM$\rangle$B]}  
	 Il tuo messaggio è stato inoltrato all'altro giocatore.\\ \tt \\ \tt \\ \tt You are currently playing round 5 of maximum 28.\\ \tt You have collectively accumulated 1 of 16 penalties.\\ \tt L'altro giocatore ha inviato questo messaggio:\\ \tt \\ \tt \\ \tt "Capito, sulla mia griglia (3,1) non è una cella libera. Aggiorno piano: propongo di usare la colonna x=4 come bersaglio comune invece di x=3. Ordine dall’alto al basso A,D,M,O,P,S,T su x=4, y1–y7. Confermi? Se sì, inizierò mettendo A su (4,1) se libero; altrimenti indicami una colonna libera alternativa."\\ \tt \\ \tt \\ \tt Qual è il tuo prossimo comando? 
	  } 
	   } 
	   } 
	 & \\ 
 

    \theutterance \stepcounter{utterance}  

    & & & \multicolumn{4}{p{0.6\linewidth}}{\cellcolor[rgb]{1,0.85,0.72}{%
	\makecell[{{p{\linewidth}}}]{% 
	  \tt {\tiny [GM$\langle$B]}  
	 COMUNICA: Ok, confermo cambio target alla colonna x=4 con ordine A,D,M,O,P,S,T dall’alto (y1–y7). Sulla mia griglia x=4 è libera alle righe 1,2,4,6,7; occupate al momento: y3 da O (posso spostarlo) e y5 da D. Procedo a liberare y1–y3 subito; quando pronto ti confermo per A@(4,1), D@(4,2), M@(4,3). 
	  } 
	   } 
	   } 
	  \\ 
 

    \theutterance \stepcounter{utterance}  

    & & & \multicolumn{2}{p{0.3\linewidth}}{\cellcolor[rgb]{0.95,0.95,0.95}{%
	\makecell[{{p{\linewidth}}}]{% 
	  \tt {\tiny [GM$|$GM]}  
	 COMUNICA: Ok, confermo cambio target alla colonna x=4 con ordine A,D,M,O,P,S,T dall’alto (y1–y7). Sulla mia griglia x=4 è libera alle righe 1,2,4,6,7; occupate al momento: y3 da O (posso spostarlo) e y5 da D. Procedo a liberare y1–y3 subito; quando pronto ti confermo per A@(4,1), D@(4,2), M@(4,3). 
	  } 
	   } 
	   } 
	 & & \\ 
 

    \theutterance \stepcounter{utterance}  

    & & \multicolumn{4}{p{0.6\linewidth}}{\cellcolor[rgb]{0.9,0.9,0.9}{%
	\makecell[{{p{\linewidth}}}]{% 
	  \tt {\tiny [A$\langle$GM]}  
	 Il tuo messaggio è stato inoltrato all'altro giocatore.\\ \tt \\ \tt \\ \tt You are currently playing round 5 of maximum 28.\\ \tt You have collectively accumulated 1 of 16 penalties.\\ \tt L'altro giocatore ha inviato questo messaggio:\\ \tt \\ \tt \\ \tt "Ok, confermo cambio target alla colonna x=4 con ordine A,D,M,O,P,S,T dall’alto (y1–y7). Sulla mia griglia x=4 è libera alle righe 1,2,4,6,7; occupate al momento: y3 da O (posso spostarlo) e y5 da D. Procedo a liberare y1–y3 subito; quando pronto ti confermo per A@(4,1), D@(4,2), M@(4,3)."\\ \tt \\ \tt \\ \tt Qual è il tuo prossimo comando? 
	  } 
	   } 
	   } 
	 & \\ 
 

    \theutterance \stepcounter{utterance}  

    & \multicolumn{4}{p{0.6\linewidth}}{\cellcolor[rgb]{0.8,1,0.9}{%
	\makecell[{{p{\linewidth}}}]{% 
	  \tt {\tiny [A$\rangle$GM]}  
	 COMUNICA: Perfetto, attendo conferma quando x=4,y1–y3 sono libere. Appena dai il via, sposto A@(4,1), poi D@(4,2), poi M@(4,3), confermando ogni passo. 
	  } 
	   } 
	   } 
	 & & \\ 
 

    \theutterance \stepcounter{utterance}  

    & & & \multicolumn{2}{p{0.3\linewidth}}{\cellcolor[rgb]{0.95,0.95,0.95}{%
	\makecell[{{p{\linewidth}}}]{% 
	  \tt {\tiny [GM$|$GM]}  
	 COMUNICA: Perfetto, attendo conferma quando x=4,y1–y3 sono libere. Appena dai il via, sposto A@(4,1), poi D@(4,2), poi M@(4,3), confermando ogni passo. 
	  } 
	   } 
	   } 
	 & & \\ 
 

    \theutterance \stepcounter{utterance}  

    & & \multicolumn{4}{p{0.6\linewidth}}{\cellcolor[rgb]{0.9,0.9,0.9}{%
	\makecell[{{p{\linewidth}}}]{% 
	  \tt {\tiny [GM$\rangle$B]}  
	 Il tuo messaggio è stato inoltrato all'altro giocatore.\\ \tt \\ \tt \\ \tt You are currently playing round 6 of maximum 28.\\ \tt You have collectively accumulated 1 of 16 penalties.\\ \tt L'altro giocatore ha inviato questo messaggio:\\ \tt \\ \tt \\ \tt "Perfetto, attendo conferma quando x=4,y1–y3 sono libere. Appena dai il via, sposto A@(4,1), poi D@(4,2), poi M@(4,3), confermando ogni passo."\\ \tt \\ \tt \\ \tt Qual è il tuo prossimo comando? 
	  } 
	   } 
	   } 
	 & \\ 
 

    \theutterance \stepcounter{utterance}  

    & & & \multicolumn{4}{p{0.6\linewidth}}{\cellcolor[rgb]{1,0.85,0.72}{%
	\makecell[{{p{\linewidth}}}]{% 
	  \tt {\tiny [GM$\langle$B]}  
	 SPOSTA: O, (3,3) 
	  } 
	   } 
	   } 
	  \\ 
 

    \theutterance \stepcounter{utterance}  

    & & & \multicolumn{2}{p{0.3\linewidth}}{\cellcolor[rgb]{0.95,0.95,0.95}{%
	\makecell[{{p{\linewidth}}}]{% 
	  \tt {\tiny [GM$|$GM]}  
	 SPOSTA: O, (3,3) 
	  } 
	   } 
	   } 
	 & & \\ 
 

    \theutterance \stepcounter{utterance}  

    & & & \multicolumn{2}{p{0.3\linewidth}}{\cellcolor[rgb]{0.95,0.95,0.95}{%
	\makecell[{{p{\linewidth}}}]{% 
	  \tt {\tiny [GM$|$GM]}  
	 Spostato 'O' in (3,3) con successo. La tua griglia aggiornata appare così:\\ \tt ```\\ \tt  1234567\\ \tt ╔══╤════╗\\ \tt ║◌S│◌◌◌◌║ 1\\ \tt ╟──┘◌◌┌─╢ 2\\ \tt ║A◌O◌◌│◌║ 3\\ \tt ╟─┐◌┌─┼─╢ 4\\ \tt ║◌├─┤T│D║ 5\\ \tt ╟─┤◌├─┤◌║ 6\\ \tt ║P│◌│M│◌║ 7\\ \tt ╚═╧═╧═╧═╝\\ \tt ```\\ \tt  
	  } 
	   } 
	   } 
	 & & \\ 
 

    \theutterance \stepcounter{utterance}  

    & & \multicolumn{4}{p{0.6\linewidth}}{\cellcolor[rgb]{0.9,0.9,0.9}{%
	\makecell[{{p{\linewidth}}}]{% 
	  \tt {\tiny [A$\langle$GM]}  
	 Il tuo messaggio è stato inoltrato all'altro giocatore.\\ \tt \\ \tt \\ \tt You are currently playing round 6 of maximum 28.\\ \tt You have collectively accumulated 1 of 16 penalties.\\ \tt L'altro giocatore ha spostato un oggetto sulla sua griglia.\\ \tt \\ \tt \\ \tt Qual è il tuo prossimo comando? 
	  } 
	   } 
	   } 
	 & \\ 
 

    \theutterance \stepcounter{utterance}  

    & \multicolumn{4}{p{0.6\linewidth}}{\cellcolor[rgb]{0.8,1,0.9}{%
	\makecell[{{p{\linewidth}}}]{% 
	  \tt {\tiny [A$\rangle$GM]}  
	 COMUNICA: Dimmi quando x=4,y1–y3 sono libere così procedo con A@(4,1), poi D@(4,2), poi M@(4,3). Pronto a muovere appena confermi. 
	  } 
	   } 
	   } 
	 & & \\ 
 

    \theutterance \stepcounter{utterance}  

    & & & \multicolumn{2}{p{0.3\linewidth}}{\cellcolor[rgb]{0.95,0.95,0.95}{%
	\makecell[{{p{\linewidth}}}]{% 
	  \tt {\tiny [GM$|$GM]}  
	 COMUNICA: Dimmi quando x=4,y1–y3 sono libere così procedo con A@(4,1), poi D@(4,2), poi M@(4,3). Pronto a muovere appena confermi. 
	  } 
	   } 
	   } 
	 & & \\ 
 

    \theutterance \stepcounter{utterance}  

    & & \multicolumn{4}{p{0.6\linewidth}}{\cellcolor[rgb]{0.9,0.9,0.9}{%
	\makecell[{{p{\linewidth}}}]{% 
	  \tt {\tiny [GM$\rangle$B]}  
	 Spostato 'O' in (3,3) con successo. La tua griglia aggiornata appare così:\\ \tt ```\\ \tt  1234567\\ \tt ╔══╤════╗\\ \tt ║◌S│◌◌◌◌║ 1\\ \tt ╟──┘◌◌┌─╢ 2\\ \tt ║A◌O◌◌│◌║ 3\\ \tt ╟─┐◌┌─┼─╢ 4\\ \tt ║◌├─┤T│D║ 5\\ \tt ╟─┤◌├─┤◌║ 6\\ \tt ║P│◌│M│◌║ 7\\ \tt ╚═╧═╧═╧═╝\\ \tt ```\\ \tt You are currently playing round 7 of maximum 28.\\ \tt You have collectively accumulated 1 of 16 penalties.\\ \tt L'altro giocatore ha inviato questo messaggio:\\ \tt \\ \tt \\ \tt "Dimmi quando x=4,y1–y3 sono libere così procedo con A@(4,1), poi D@(4,2), poi M@(4,3). Pronto a muovere appena confermi."\\ \tt \\ \tt \\ \tt Qual è il tuo prossimo comando? 
	  } 
	   } 
	   } 
	 & \\ 
 

    \theutterance \stepcounter{utterance}  

    & & & \multicolumn{4}{p{0.6\linewidth}}{\cellcolor[rgb]{1,0.85,0.72}{%
	\makecell[{{p{\linewidth}}}]{% 
	  \tt {\tiny [GM$\langle$B]}  
	 COMUNICA: x=4,y1–y3 sono libere ora. Puoi muovere in ordine: A@(4,1), poi D@(4,2), poi M@(4,3). Conferma ogni passo e poi preparo O@(4,4), P@(4,5), S@(4,6), T@(4,7). 
	  } 
	   } 
	   } 
	  \\ 
 

    \theutterance \stepcounter{utterance}  

    & & & \multicolumn{2}{p{0.3\linewidth}}{\cellcolor[rgb]{0.95,0.95,0.95}{%
	\makecell[{{p{\linewidth}}}]{% 
	  \tt {\tiny [GM$|$GM]}  
	 COMUNICA: x=4,y1–y3 sono libere ora. Puoi muovere in ordine: A@(4,1), poi D@(4,2), poi M@(4,3). Conferma ogni passo e poi preparo O@(4,4), P@(4,5), S@(4,6), T@(4,7). 
	  } 
	   } 
	   } 
	 & & \\ 
 

    \theutterance \stepcounter{utterance}  

    & & \multicolumn{4}{p{0.6\linewidth}}{\cellcolor[rgb]{0.9,0.9,0.9}{%
	\makecell[{{p{\linewidth}}}]{% 
	  \tt {\tiny [A$\langle$GM]}  
	 Il tuo messaggio è stato inoltrato all'altro giocatore.\\ \tt \\ \tt \\ \tt You are currently playing round 7 of maximum 28.\\ \tt You have collectively accumulated 1 of 16 penalties.\\ \tt L'altro giocatore ha inviato questo messaggio:\\ \tt \\ \tt \\ \tt "x=4,y1–y3 sono libere ora. Puoi muovere in ordine: A@(4,1), poi D@(4,2), poi M@(4,3). Conferma ogni passo e poi preparo O@(4,4), P@(4,5), S@(4,6), T@(4,7)."\\ \tt \\ \tt \\ \tt Qual è il tuo prossimo comando? 
	  } 
	   } 
	   } 
	 & \\ 
 

    \theutterance \stepcounter{utterance}  

    & \multicolumn{4}{p{0.6\linewidth}}{\cellcolor[rgb]{0.8,1,0.9}{%
	\makecell[{{p{\linewidth}}}]{% 
	  \tt {\tiny [A$\rangle$GM]}  
	 SPOSTA: A, (4, 1) 
	  } 
	   } 
	   } 
	 & & \\ 
 

    \theutterance \stepcounter{utterance}  

    & & & \multicolumn{2}{p{0.3\linewidth}}{\cellcolor[rgb]{0.95,0.95,0.95}{%
	\makecell[{{p{\linewidth}}}]{% 
	  \tt {\tiny [GM$|$GM]}  
	 SPOSTA: A, (4, 1) 
	  } 
	   } 
	   } 
	 & & \\ 
 

    \theutterance \stepcounter{utterance}  

    & & & \multicolumn{2}{p{0.3\linewidth}}{\cellcolor[rgb]{0.95,0.95,0.95}{%
	\makecell[{{p{\linewidth}}}]{% 
	  \tt {\tiny [GM$|$GM]}  
	 Penalità: (4,1) non è vuoto, ma contiene 'D'.\\ \tt \\ \tt You have collectively accumulated 2 of 16 penalties.\\ \tt Riprova! 
	  } 
	   } 
	   } 
	 & & \\ 
 

    \theutterance \stepcounter{utterance}  

    & & \multicolumn{4}{p{0.6\linewidth}}{\cellcolor[rgb]{0.9,0.9,0.9}{%
	\makecell[{{p{\linewidth}}}]{% 
	  \tt {\tiny [A$\langle$GM]}  
	 Penalità: (4,1) non è vuoto, ma contiene 'D'.\\ \tt \\ \tt You have collectively accumulated 2 of 16 penalties.\\ \tt Riprova! 
	  } 
	   } 
	   } 
	 & \\ 
 

    \theutterance \stepcounter{utterance}  

    & \multicolumn{4}{p{0.6\linewidth}}{\cellcolor[rgb]{0.8,1,0.9}{%
	\makecell[{{p{\linewidth}}}]{% 
	  \tt {\tiny [A$\rangle$GM]}  
	 COMUNICA: Sulla mia griglia (4,1) è occupato da D. Propongo di riempire la colonna x=4 dall’alto verso il basso saltando le celle occupate: posiziono A nella prima riga libera partendo dall’alto. Verifico: (4,2) è libero? Se sì, muovo A@(4,2), poi D@(4,3) e M@(4,4), mantenendo l’ordine relativo A,D,M dall’alto verso il basso e lasciando a te O,P,S,T sotto. Conferma quali righe libere hai su x=4 nelle posizioni y1–y4 dal tuo lato, così restiamo sincronizzati. 
	  } 
	   } 
	   } 
	 & & \\ 
 

    \theutterance \stepcounter{utterance}  

    & & & \multicolumn{2}{p{0.3\linewidth}}{\cellcolor[rgb]{0.95,0.95,0.95}{%
	\makecell[{{p{\linewidth}}}]{% 
	  \tt {\tiny [GM$|$GM]}  
	 COMUNICA: Sulla mia griglia (4,1) è occupato da D. Propongo di riempire la colonna x=4 dall’alto verso il basso saltando le celle occupate: posiziono A nella prima riga libera partendo dall’alto. Verifico: (4,2) è libero? Se sì, muovo A@(4,2), poi D@(4,3) e M@(4,4), mantenendo l’ordine relativo A,D,M dall’alto verso il basso e lasciando a te O,P,S,T sotto. Conferma quali righe libere hai su x=4 nelle posizioni y1–y4 dal tuo lato, così restiamo sincronizzati. 
	  } 
	   } 
	   } 
	 & & \\ 
 

    \theutterance \stepcounter{utterance}  

    & & \multicolumn{4}{p{0.6\linewidth}}{\cellcolor[rgb]{0.9,0.9,0.9}{%
	\makecell[{{p{\linewidth}}}]{% 
	  \tt {\tiny [GM$\rangle$B]}  
	 Il tuo messaggio è stato inoltrato all'altro giocatore.\\ \tt \\ \tt \\ \tt You are currently playing round 8 of maximum 28.\\ \tt You have collectively accumulated 2 of 16 penalties.\\ \tt L'altro giocatore ha inviato questo messaggio:\\ \tt \\ \tt \\ \tt "Sulla mia griglia (4,1) è occupato da D. Propongo di riempire la colonna x=4 dall’alto verso il basso saltando le celle occupate: posiziono A nella prima riga libera partendo dall’alto. Verifico: (4,2) è libero? Se sì, muovo A@(4,2), poi D@(4,3) e M@(4,4), mantenendo l’ordine relativo A,D,M dall’alto verso il basso e lasciando a te O,P,S,T sotto. Conferma quali righe libere hai su x=4 nelle posizioni y1–y4 dal tuo lato, così restiamo sincronizzati."\\ \tt \\ \tt \\ \tt Qual è il tuo prossimo comando? 
	  } 
	   } 
	   } 
	 & \\ 
 

    \theutterance \stepcounter{utterance}  

    & & & \multicolumn{4}{p{0.6\linewidth}}{\cellcolor[rgb]{1,0.85,0.72}{%
	\makecell[{{p{\linewidth}}}]{% 
	  \tt {\tiny [GM$\langle$B]}  
	 COMUNICA: Sul mio lato x=4 è così: y1 libero, y2 libero, y3 libero, y4 libero, y5 occupato da D, y6 libero, y7 libero. Per sincronizzarci con il tuo vincolo, confermo che puoi posizionare A nella prima riga libera partendo dall’alto sul tuo lato (se (4,1) è occupato, usa (4,2)), poi D sotto A, poi M sotto D. Io manterrò l’ordine A,D,M in cima anche da me: appena mi confermi le tue posizioni effettive, passerò a O@(4,4), P@(4,5), S@(4,6), T@(4,7) per restare allineati. 
	  } 
	   } 
	   } 
	  \\ 
 

    \theutterance \stepcounter{utterance}  

    & & & \multicolumn{2}{p{0.3\linewidth}}{\cellcolor[rgb]{0.95,0.95,0.95}{%
	\makecell[{{p{\linewidth}}}]{% 
	  \tt {\tiny [GM$|$GM]}  
	 COMUNICA: Sul mio lato x=4 è così: y1 libero, y2 libero, y3 libero, y4 libero, y5 occupato da D, y6 libero, y7 libero. Per sincronizzarci con il tuo vincolo, confermo che puoi posizionare A nella prima riga libera partendo dall’alto sul tuo lato (se (4,1) è occupato, usa (4,2)), poi D sotto A, poi M sotto D. Io manterrò l’ordine A,D,M in cima anche da me: appena mi confermi le tue posizioni effettive, passerò a O@(4,4), P@(4,5), S@(4,6), T@(4,7) per restare allineati. 
	  } 
	   } 
	   } 
	 & & \\ 
 

    \theutterance \stepcounter{utterance}  

    & & \multicolumn{4}{p{0.6\linewidth}}{\cellcolor[rgb]{0.9,0.9,0.9}{%
	\makecell[{{p{\linewidth}}}]{% 
	  \tt {\tiny [A$\langle$GM]}  
	 Il tuo messaggio è stato inoltrato all'altro giocatore.\\ \tt \\ \tt \\ \tt You are currently playing round 8 of maximum 28.\\ \tt You have collectively accumulated 2 of 16 penalties.\\ \tt L'altro giocatore ha inviato questo messaggio:\\ \tt \\ \tt \\ \tt "Sul mio lato x=4 è così: y1 libero, y2 libero, y3 libero, y4 libero, y5 occupato da D, y6 libero, y7 libero. Per sincronizzarci con il tuo vincolo, confermo che puoi posizionare A nella prima riga libera partendo dall’alto sul tuo lato (se (4,1) è occupato, usa (4,2)), poi D sotto A, poi M sotto D. Io manterrò l’ordine A,D,M in cima anche da me: appena mi confermi le tue posizioni effettive, passerò a O@(4,4), P@(4,5), S@(4,6), T@(4,7) per restare allineati."\\ \tt \\ \tt \\ \tt Qual è il tuo prossimo comando? 
	  } 
	   } 
	   } 
	 & \\ 
 

    \theutterance \stepcounter{utterance}  

    & \multicolumn{4}{p{0.6\linewidth}}{\cellcolor[rgb]{0.8,1,0.9}{%
	\makecell[{{p{\linewidth}}}]{% 
	  \tt {\tiny [A$\rangle$GM]}  
	 SPOSTA: A, (4, 2) 
	  } 
	   } 
	   } 
	 & & \\ 
 

    \theutterance \stepcounter{utterance}  

    & & & \multicolumn{2}{p{0.3\linewidth}}{\cellcolor[rgb]{0.95,0.95,0.95}{%
	\makecell[{{p{\linewidth}}}]{% 
	  \tt {\tiny [GM$|$GM]}  
	 SPOSTA: A, (4, 2) 
	  } 
	   } 
	   } 
	 & & \\ 
 

    \theutterance \stepcounter{utterance}  

    & & & \multicolumn{2}{p{0.3\linewidth}}{\cellcolor[rgb]{0.95,0.95,0.95}{%
	\makecell[{{p{\linewidth}}}]{% 
	  \tt {\tiny [GM$|$GM]}  
	 Spostato 'A' in (4,2) con successo. La tua griglia aggiornata appare così:\\ \tt ```\\ \tt  1234567\\ \tt ╔══╤════╗\\ \tt ║◌T│D◌◌◌║ 1\\ \tt ╟──┘A◌┌─╢ 2\\ \tt ║S◌◌◌◌│◌║ 3\\ \tt ╟─┐◌┌─┼─╢ 4\\ \tt ║◌├─┤◌│O║ 5\\ \tt ╟─┤◌├─┤◌║ 6\\ \tt ║◌│P│M│◌║ 7\\ \tt ╚═╧═╧═╧═╝\\ \tt ```\\ \tt  
	  } 
	   } 
	   } 
	 & & \\ 
 

    \theutterance \stepcounter{utterance}  

    & & \multicolumn{4}{p{0.6\linewidth}}{\cellcolor[rgb]{0.9,0.9,0.9}{%
	\makecell[{{p{\linewidth}}}]{% 
	  \tt {\tiny [GM$\rangle$B]}  
	 Il tuo messaggio è stato inoltrato all'altro giocatore.\\ \tt \\ \tt \\ \tt You are currently playing round 9 of maximum 28.\\ \tt You have collectively accumulated 2 of 16 penalties.\\ \tt L'altro giocatore ha spostato un oggetto sulla sua griglia.\\ \tt \\ \tt \\ \tt Qual è il tuo prossimo comando? 
	  } 
	   } 
	   } 
	 & \\ 
 

    \theutterance \stepcounter{utterance}  

    & & & \multicolumn{4}{p{0.6\linewidth}}{\cellcolor[rgb]{1,0.85,0.72}{%
	\makecell[{{p{\linewidth}}}]{% 
	  \tt {\tiny [GM$\langle$B]}  
	 SPOSTA: D, (4,5) 
	  } 
	   } 
	   } 
	  \\ 
 

    \theutterance \stepcounter{utterance}  

    & & & \multicolumn{2}{p{0.3\linewidth}}{\cellcolor[rgb]{0.95,0.95,0.95}{%
	\makecell[{{p{\linewidth}}}]{% 
	  \tt {\tiny [GM$|$GM]}  
	 SPOSTA: D, (4,5) 
	  } 
	   } 
	   } 
	 & & \\ 
 

    \theutterance \stepcounter{utterance}  

    & & & \multicolumn{2}{p{0.3\linewidth}}{\cellcolor[rgb]{0.95,0.95,0.95}{%
	\makecell[{{p{\linewidth}}}]{% 
	  \tt {\tiny [GM$|$GM]}  
	 Penalità: (4,5) non è vuoto, ma contiene '┤'.\\ \tt \\ \tt You have collectively accumulated 3 of 16 penalties.\\ \tt Riprova! 
	  } 
	   } 
	   } 
	 & & \\ 
 

    \theutterance \stepcounter{utterance}  

    & & \multicolumn{4}{p{0.6\linewidth}}{\cellcolor[rgb]{0.9,0.9,0.9}{%
	\makecell[{{p{\linewidth}}}]{% 
	  \tt {\tiny [GM$\rangle$B]}  
	 Penalità: (4,5) non è vuoto, ma contiene '┤'.\\ \tt \\ \tt You have collectively accumulated 3 of 16 penalties.\\ \tt Riprova! 
	  } 
	   } 
	   } 
	 & \\ 
 

    \theutterance \stepcounter{utterance}  

    & & & \multicolumn{4}{p{0.6\linewidth}}{\cellcolor[rgb]{1,0.85,0.72}{%
	\makecell[{{p{\linewidth}}}]{% 
	  \tt {\tiny [GM$\langle$B]}  
	 SPOSTA: D, (4,2) 
	  } 
	   } 
	   } 
	  \\ 
 

    \theutterance \stepcounter{utterance}  

    & & & \multicolumn{2}{p{0.3\linewidth}}{\cellcolor[rgb]{0.95,0.95,0.95}{%
	\makecell[{{p{\linewidth}}}]{% 
	  \tt {\tiny [GM$|$GM]}  
	 SPOSTA: D, (4,2) 
	  } 
	   } 
	   } 
	 & & \\ 
 

    \theutterance \stepcounter{utterance}  

    & & & \multicolumn{2}{p{0.3\linewidth}}{\cellcolor[rgb]{0.95,0.95,0.95}{%
	\makecell[{{p{\linewidth}}}]{% 
	  \tt {\tiny [GM$|$GM]}  
	 Spostato 'D' in (4,2) con successo. La tua griglia aggiornata appare così:\\ \tt ```\\ \tt  1234567\\ \tt ╔══╤════╗\\ \tt ║◌S│◌◌◌◌║ 1\\ \tt ╟──┘D◌┌─╢ 2\\ \tt ║A◌O◌◌│◌║ 3\\ \tt ╟─┐◌┌─┼─╢ 4\\ \tt ║◌├─┤T│◌║ 5\\ \tt ╟─┤◌├─┤◌║ 6\\ \tt ║P│◌│M│◌║ 7\\ \tt ╚═╧═╧═╧═╝\\ \tt ```\\ \tt  
	  } 
	   } 
	   } 
	 & & \\ 
 

    \theutterance \stepcounter{utterance}  

    & & \multicolumn{4}{p{0.6\linewidth}}{\cellcolor[rgb]{0.9,0.9,0.9}{%
	\makecell[{{p{\linewidth}}}]{% 
	  \tt {\tiny [A$\langle$GM]}  
	 Spostato 'A' in (4,2) con successo. La tua griglia aggiornata appare così:\\ \tt ```\\ \tt  1234567\\ \tt ╔══╤════╗\\ \tt ║◌T│D◌◌◌║ 1\\ \tt ╟──┘A◌┌─╢ 2\\ \tt ║S◌◌◌◌│◌║ 3\\ \tt ╟─┐◌┌─┼─╢ 4\\ \tt ║◌├─┤◌│O║ 5\\ \tt ╟─┤◌├─┤◌║ 6\\ \tt ║◌│P│M│◌║ 7\\ \tt ╚═╧═╧═╧═╝\\ \tt ```\\ \tt You are currently playing round 9 of maximum 28.\\ \tt You have collectively accumulated 3 of 16 penalties.\\ \tt L'altro giocatore ha spostato un oggetto sulla sua griglia.\\ \tt \\ \tt \\ \tt Qual è il tuo prossimo comando? 
	  } 
	   } 
	   } 
	 & \\ 
 

    \theutterance \stepcounter{utterance}  

    & \multicolumn{4}{p{0.6\linewidth}}{\cellcolor[rgb]{0.8,1,0.9}{%
	\makecell[{{p{\linewidth}}}]{% 
	  \tt {\tiny [A$\rangle$GM]}  
	 SPOSTA: D, (4, 3) 
	  } 
	   } 
	   } 
	 & & \\ 
 

    \theutterance \stepcounter{utterance}  

    & & & \multicolumn{2}{p{0.3\linewidth}}{\cellcolor[rgb]{0.95,0.95,0.95}{%
	\makecell[{{p{\linewidth}}}]{% 
	  \tt {\tiny [GM$|$GM]}  
	 SPOSTA: D, (4, 3) 
	  } 
	   } 
	   } 
	 & & \\ 
 

    \theutterance \stepcounter{utterance}  

    & & & \multicolumn{2}{p{0.3\linewidth}}{\cellcolor[rgb]{0.95,0.95,0.95}{%
	\makecell[{{p{\linewidth}}}]{% 
	  \tt {\tiny [GM$|$GM]}  
	 Spostato 'D' in (4,3) con successo. La tua griglia aggiornata appare così:\\ \tt ```\\ \tt  1234567\\ \tt ╔══╤════╗\\ \tt ║◌T│◌◌◌◌║ 1\\ \tt ╟──┘A◌┌─╢ 2\\ \tt ║S◌◌D◌│◌║ 3\\ \tt ╟─┐◌┌─┼─╢ 4\\ \tt ║◌├─┤◌│O║ 5\\ \tt ╟─┤◌├─┤◌║ 6\\ \tt ║◌│P│M│◌║ 7\\ \tt ╚═╧═╧═╧═╝\\ \tt ```\\ \tt  
	  } 
	   } 
	   } 
	 & & \\ 
 

    \theutterance \stepcounter{utterance}  

    & & \multicolumn{4}{p{0.6\linewidth}}{\cellcolor[rgb]{0.9,0.9,0.9}{%
	\makecell[{{p{\linewidth}}}]{% 
	  \tt {\tiny [GM$\rangle$B]}  
	 Spostato 'D' in (4,2) con successo. La tua griglia aggiornata appare così:\\ \tt ```\\ \tt  1234567\\ \tt ╔══╤════╗\\ \tt ║◌S│◌◌◌◌║ 1\\ \tt ╟──┘D◌┌─╢ 2\\ \tt ║A◌O◌◌│◌║ 3\\ \tt ╟─┐◌┌─┼─╢ 4\\ \tt ║◌├─┤T│◌║ 5\\ \tt ╟─┤◌├─┤◌║ 6\\ \tt ║P│◌│M│◌║ 7\\ \tt ╚═╧═╧═╧═╝\\ \tt ```\\ \tt You are currently playing round 10 of maximum 28.\\ \tt You have collectively accumulated 3 of 16 penalties.\\ \tt L'altro giocatore ha spostato un oggetto sulla sua griglia.\\ \tt \\ \tt \\ \tt Qual è il tuo prossimo comando? 
	  } 
	   } 
	   } 
	 & \\ 
 

    \theutterance \stepcounter{utterance}  

    & & & \multicolumn{4}{p{0.6\linewidth}}{\cellcolor[rgb]{1,0.85,0.72}{%
	\makecell[{{p{\linewidth}}}]{% 
	  \tt {\tiny [GM$\langle$B]}  
	 COMUNICA: Ho spostato D su (4,2). Stato colonna x=4 ora: y1 libero, y2=D, y3 libero, y4 libero, y5 libero, y6 libero, y7 libero. Conferma dove hai messo A e se procedi con D sotto A e M sotto D. Appena confermi, io metto O su (4,3) per mantenere l’ordine A,D,M,O,P,S,T dall’alto verso il basso. 
	  } 
	   } 
	   } 
	  \\ 
 

    \theutterance \stepcounter{utterance}  

    & & & \multicolumn{2}{p{0.3\linewidth}}{\cellcolor[rgb]{0.95,0.95,0.95}{%
	\makecell[{{p{\linewidth}}}]{% 
	  \tt {\tiny [GM$|$GM]}  
	 COMUNICA: Ho spostato D su (4,2). Stato colonna x=4 ora: y1 libero, y2=D, y3 libero, y4 libero, y5 libero, y6 libero, y7 libero. Conferma dove hai messo A e se procedi con D sotto A e M sotto D. Appena confermi, io metto O su (4,3) per mantenere l’ordine A,D,M,O,P,S,T dall’alto verso il basso. 
	  } 
	   } 
	   } 
	 & & \\ 
 

    \theutterance \stepcounter{utterance}  

    & & \multicolumn{4}{p{0.6\linewidth}}{\cellcolor[rgb]{0.9,0.9,0.9}{%
	\makecell[{{p{\linewidth}}}]{% 
	  \tt {\tiny [A$\langle$GM]}  
	 Spostato 'D' in (4,3) con successo. La tua griglia aggiornata appare così:\\ \tt ```\\ \tt  1234567\\ \tt ╔══╤════╗\\ \tt ║◌T│◌◌◌◌║ 1\\ \tt ╟──┘A◌┌─╢ 2\\ \tt ║S◌◌D◌│◌║ 3\\ \tt ╟─┐◌┌─┼─╢ 4\\ \tt ║◌├─┤◌│O║ 5\\ \tt ╟─┤◌├─┤◌║ 6\\ \tt ║◌│P│M│◌║ 7\\ \tt ╚═╧═╧═╧═╝\\ \tt ```\\ \tt You are currently playing round 10 of maximum 28.\\ \tt You have collectively accumulated 3 of 16 penalties.\\ \tt L'altro giocatore ha inviato questo messaggio:\\ \tt \\ \tt \\ \tt "Ho spostato D su (4,2). Stato colonna x=4 ora: y1 libero, y2=D, y3 libero, y4 libero, y5 libero, y6 libero, y7 libero. Conferma dove hai messo A e se procedi con D sotto A e M sotto D. Appena confermi, io metto O su (4,3) per mantenere l’ordine A,D,M,O,P,S,T dall’alto verso il basso."\\ \tt \\ \tt \\ \tt Qual è il tuo prossimo comando? 
	  } 
	   } 
	   } 
	 & \\ 
 

    \theutterance \stepcounter{utterance}  

    & \multicolumn{4}{p{0.6\linewidth}}{\cellcolor[rgb]{0.8,1,0.9}{%
	\makecell[{{p{\linewidth}}}]{% 
	  \tt {\tiny [A$\rangle$GM]}  
	 COMUNICA: Confermo: A@(4,2), D@(4,3) posizionato, prossimo M@(4,4). Dopo metto O@(4,5), P@(4,6), S@(4,7) se liberi; tu puoi mettere O@(4,4) sul tuo lato per mantenere A,D,M,O dall’alto. Ti avviso appena M è su (4,4). 
	  } 
	   } 
	   } 
	 & & \\ 
 

    \theutterance \stepcounter{utterance}  

    & & & \multicolumn{2}{p{0.3\linewidth}}{\cellcolor[rgb]{0.95,0.95,0.95}{%
	\makecell[{{p{\linewidth}}}]{% 
	  \tt {\tiny [GM$|$GM]}  
	 COMUNICA: Confermo: A@(4,2), D@(4,3) posizionato, prossimo M@(4,4). Dopo metto O@(4,5), P@(4,6), S@(4,7) se liberi; tu puoi mettere O@(4,4) sul tuo lato per mantenere A,D,M,O dall’alto. Ti avviso appena M è su (4,4). 
	  } 
	   } 
	   } 
	 & & \\ 
 

    \theutterance \stepcounter{utterance}  

    & & \multicolumn{4}{p{0.6\linewidth}}{\cellcolor[rgb]{0.9,0.9,0.9}{%
	\makecell[{{p{\linewidth}}}]{% 
	  \tt {\tiny [GM$\rangle$B]}  
	 Il tuo messaggio è stato inoltrato all'altro giocatore.\\ \tt \\ \tt \\ \tt You are currently playing round 11 of maximum 28.\\ \tt You have collectively accumulated 3 of 16 penalties.\\ \tt L'altro giocatore ha inviato questo messaggio:\\ \tt \\ \tt \\ \tt "Confermo: A@(4,2), D@(4,3) posizionato, prossimo M@(4,4). Dopo metto O@(4,5), P@(4,6), S@(4,7) se liberi; tu puoi mettere O@(4,4) sul tuo lato per mantenere A,D,M,O dall’alto. Ti avviso appena M è su (4,4)."\\ \tt \\ \tt \\ \tt Qual è il tuo prossimo comando? 
	  } 
	   } 
	   } 
	 & \\ 
 

    \theutterance \stepcounter{utterance}  

    & & & \multicolumn{4}{p{0.6\linewidth}}{\cellcolor[rgb]{1,0.85,0.72}{%
	\makecell[{{p{\linewidth}}}]{% 
	  \tt {\tiny [GM$\langle$B]}  
	 SPOSTA: O, (4,3) 
	  } 
	   } 
	   } 
	  \\ 
 

    \theutterance \stepcounter{utterance}  

    & & & \multicolumn{2}{p{0.3\linewidth}}{\cellcolor[rgb]{0.95,0.95,0.95}{%
	\makecell[{{p{\linewidth}}}]{% 
	  \tt {\tiny [GM$|$GM]}  
	 SPOSTA: O, (4,3) 
	  } 
	   } 
	   } 
	 & & \\ 
 

    \theutterance \stepcounter{utterance}  

    & & & \multicolumn{2}{p{0.3\linewidth}}{\cellcolor[rgb]{0.95,0.95,0.95}{%
	\makecell[{{p{\linewidth}}}]{% 
	  \tt {\tiny [GM$|$GM]}  
	 Spostato 'O' in (4,3) con successo. La tua griglia aggiornata appare così:\\ \tt ```\\ \tt  1234567\\ \tt ╔══╤════╗\\ \tt ║◌S│◌◌◌◌║ 1\\ \tt ╟──┘D◌┌─╢ 2\\ \tt ║A◌◌O◌│◌║ 3\\ \tt ╟─┐◌┌─┼─╢ 4\\ \tt ║◌├─┤T│◌║ 5\\ \tt ╟─┤◌├─┤◌║ 6\\ \tt ║P│◌│M│◌║ 7\\ \tt ╚═╧═╧═╧═╝\\ \tt ```\\ \tt  
	  } 
	   } 
	   } 
	 & & \\ 
 

    \theutterance \stepcounter{utterance}  

    & & \multicolumn{4}{p{0.6\linewidth}}{\cellcolor[rgb]{0.9,0.9,0.9}{%
	\makecell[{{p{\linewidth}}}]{% 
	  \tt {\tiny [A$\langle$GM]}  
	 Il tuo messaggio è stato inoltrato all'altro giocatore.\\ \tt \\ \tt \\ \tt You are currently playing round 11 of maximum 28.\\ \tt You have collectively accumulated 3 of 16 penalties.\\ \tt L'altro giocatore ha spostato un oggetto sulla sua griglia.\\ \tt \\ \tt \\ \tt Qual è il tuo prossimo comando? 
	  } 
	   } 
	   } 
	 & \\ 
 

    \theutterance \stepcounter{utterance}  

    & \multicolumn{4}{p{0.6\linewidth}}{\cellcolor[rgb]{0.8,1,0.9}{%
	\makecell[{{p{\linewidth}}}]{% 
	  \tt {\tiny [A$\rangle$GM]}  
	 SPOSTA: M, (4, 4) 
	  } 
	   } 
	   } 
	 & & \\ 
 

    \theutterance \stepcounter{utterance}  

    & & & \multicolumn{2}{p{0.3\linewidth}}{\cellcolor[rgb]{0.95,0.95,0.95}{%
	\makecell[{{p{\linewidth}}}]{% 
	  \tt {\tiny [GM$|$GM]}  
	 SPOSTA: M, (4, 4) 
	  } 
	   } 
	   } 
	 & & \\ 
 

    \theutterance \stepcounter{utterance}  

    & & & \multicolumn{2}{p{0.3\linewidth}}{\cellcolor[rgb]{0.95,0.95,0.95}{%
	\makecell[{{p{\linewidth}}}]{% 
	  \tt {\tiny [GM$|$GM]}  
	 Penalità: (4,4) non è vuoto, ma contiene '┌'.\\ \tt \\ \tt You have collectively accumulated 4 of 16 penalties.\\ \tt Riprova! 
	  } 
	   } 
	   } 
	 & & \\ 
 

    \theutterance \stepcounter{utterance}  

    & & \multicolumn{4}{p{0.6\linewidth}}{\cellcolor[rgb]{0.9,0.9,0.9}{%
	\makecell[{{p{\linewidth}}}]{% 
	  \tt {\tiny [A$\langle$GM]}  
	 Penalità: (4,4) non è vuoto, ma contiene '┌'.\\ \tt \\ \tt You have collectively accumulated 4 of 16 penalties.\\ \tt Riprova! 
	  } 
	   } 
	   } 
	 & \\ 
 

    \theutterance \stepcounter{utterance}  

    & \multicolumn{4}{p{0.6\linewidth}}{\cellcolor[rgb]{0.8,1,0.9}{%
	\makecell[{{p{\linewidth}}}]{% 
	  \tt {\tiny [A$\rangle$GM]}  
	 COMUNICA: Sul mio lato (4,4) è bloccato da una linea. Aggiorno piano: mantengo l’ordine A@(4,2), D@(4,3), poi cerco la prossima riga libera per M partendo da y4 verso il basso. Verifica: posso mettere M su (4,5) se è libero sul mio lato; in caso contrario userò la prima libera tra y5–y7. Tu mantieni A,D,M dall’alto e poi O,P,S,T a seguire. Conferma che per te resta ok l’ordine per righe y1–y7 indipendentemente da piccole differenze di altezza finché l’ordine A,D,M,O,P,S,T è rispettato. 
	  } 
	   } 
	   } 
	 & & \\ 
 

    \theutterance \stepcounter{utterance}  

    & & & \multicolumn{2}{p{0.3\linewidth}}{\cellcolor[rgb]{0.95,0.95,0.95}{%
	\makecell[{{p{\linewidth}}}]{% 
	  \tt {\tiny [GM$|$GM]}  
	 COMUNICA: Sul mio lato (4,4) è bloccato da una linea. Aggiorno piano: mantengo l’ordine A@(4,2), D@(4,3), poi cerco la prossima riga libera per M partendo da y4 verso il basso. Verifica: posso mettere M su (4,5) se è libero sul mio lato; in caso contrario userò la prima libera tra y5–y7. Tu mantieni A,D,M dall’alto e poi O,P,S,T a seguire. Conferma che per te resta ok l’ordine per righe y1–y7 indipendentemente da piccole differenze di altezza finché l’ordine A,D,M,O,P,S,T è rispettato. 
	  } 
	   } 
	   } 
	 & & \\ 
 

    \theutterance \stepcounter{utterance}  

    & & \multicolumn{4}{p{0.6\linewidth}}{\cellcolor[rgb]{0.9,0.9,0.9}{%
	\makecell[{{p{\linewidth}}}]{% 
	  \tt {\tiny [GM$\rangle$B]}  
	 Spostato 'O' in (4,3) con successo. La tua griglia aggiornata appare così:\\ \tt ```\\ \tt  1234567\\ \tt ╔══╤════╗\\ \tt ║◌S│◌◌◌◌║ 1\\ \tt ╟──┘D◌┌─╢ 2\\ \tt ║A◌◌O◌│◌║ 3\\ \tt ╟─┐◌┌─┼─╢ 4\\ \tt ║◌├─┤T│◌║ 5\\ \tt ╟─┤◌├─┤◌║ 6\\ \tt ║P│◌│M│◌║ 7\\ \tt ╚═╧═╧═╧═╝\\ \tt ```\\ \tt You are currently playing round 12 of maximum 28.\\ \tt You have collectively accumulated 4 of 16 penalties.\\ \tt L'altro giocatore ha inviato questo messaggio:\\ \tt \\ \tt \\ \tt "Sul mio lato (4,4) è bloccato da una linea. Aggiorno piano: mantengo l’ordine A@(4,2), D@(4,3), poi cerco la prossima riga libera per M partendo da y4 verso il basso. Verifica: posso mettere M su (4,5) se è libero sul mio lato; in caso contrario userò la prima libera tra y5–y7. Tu mantieni A,D,M dall’alto e poi O,P,S,T a seguire. Conferma che per te resta ok l’ordine per righe y1–y7 indipendentemente da piccole differenze di altezza finché l’ordine A,D,M,O,P,S,T è rispettato."\\ \tt \\ \tt \\ \tt Qual è il tuo prossimo comando? 
	  } 
	   } 
	   } 
	 & \\ 
 

    \theutterance \stepcounter{utterance}  

    & & & \multicolumn{4}{p{0.6\linewidth}}{\cellcolor[rgb]{1,0.85,0.72}{%
	\makecell[{{p{\linewidth}}}]{% 
	  \tt {\tiny [GM$\langle$B]}  
	 COMUNICA: Va bene, confermo che l’obiettivo condiviso è mantenere l’ordine A,D,M,O,P,S,T in colonna x=4 dall’alto verso il basso, anche se le altezze precise possono variare a causa di celle bloccate. Stato mio: x=4 ha D@(4,2) e O@(4,3); le altre righe sono libere. Quando confermi la posizione finale di M (prima libera da y4 in giù sul tuo lato), procederò a posizionare P@(4,4) poi S@(4,5) poi T@(4,6) sul mio, lasciando eventualmente spazio per allineare O,P,S,T sotto al tuo M. 
	  } 
	   } 
	   } 
	  \\ 
 

    \theutterance \stepcounter{utterance}  

    & & & \multicolumn{2}{p{0.3\linewidth}}{\cellcolor[rgb]{0.95,0.95,0.95}{%
	\makecell[{{p{\linewidth}}}]{% 
	  \tt {\tiny [GM$|$GM]}  
	 COMUNICA: Va bene, confermo che l’obiettivo condiviso è mantenere l’ordine A,D,M,O,P,S,T in colonna x=4 dall’alto verso il basso, anche se le altezze precise possono variare a causa di celle bloccate. Stato mio: x=4 ha D@(4,2) e O@(4,3); le altre righe sono libere. Quando confermi la posizione finale di M (prima libera da y4 in giù sul tuo lato), procederò a posizionare P@(4,4) poi S@(4,5) poi T@(4,6) sul mio, lasciando eventualmente spazio per allineare O,P,S,T sotto al tuo M. 
	  } 
	   } 
	   } 
	 & & \\ 
 

    \theutterance \stepcounter{utterance}  

    & & \multicolumn{4}{p{0.6\linewidth}}{\cellcolor[rgb]{0.9,0.9,0.9}{%
	\makecell[{{p{\linewidth}}}]{% 
	  \tt {\tiny [A$\langle$GM]}  
	 Il tuo messaggio è stato inoltrato all'altro giocatore.\\ \tt \\ \tt \\ \tt You are currently playing round 12 of maximum 28.\\ \tt You have collectively accumulated 4 of 16 penalties.\\ \tt L'altro giocatore ha inviato questo messaggio:\\ \tt \\ \tt \\ \tt "Va bene, confermo che l’obiettivo condiviso è mantenere l’ordine A,D,M,O,P,S,T in colonna x=4 dall’alto verso il basso, anche se le altezze precise possono variare a causa di celle bloccate. Stato mio: x=4 ha D@(4,2) e O@(4,3); le altre righe sono libere. Quando confermi la posizione finale di M (prima libera da y4 in giù sul tuo lato), procederò a posizionare P@(4,4) poi S@(4,5) poi T@(4,6) sul mio, lasciando eventualmente spazio per allineare O,P,S,T sotto al tuo M."\\ \tt \\ \tt \\ \tt Qual è il tuo prossimo comando? 
	  } 
	   } 
	   } 
	 & \\ 
 

    \theutterance \stepcounter{utterance}  

    & \multicolumn{4}{p{0.6\linewidth}}{\cellcolor[rgb]{0.8,1,0.9}{%
	\makecell[{{p{\linewidth}}}]{% 
	  \tt {\tiny [A$\rangle$GM]}  
	 SPOSTA: M, (4, 5) 
	  } 
	   } 
	   } 
	 & & \\ 
 

    \theutterance \stepcounter{utterance}  

    & & & \multicolumn{2}{p{0.3\linewidth}}{\cellcolor[rgb]{0.95,0.95,0.95}{%
	\makecell[{{p{\linewidth}}}]{% 
	  \tt {\tiny [GM$|$GM]}  
	 SPOSTA: M, (4, 5) 
	  } 
	   } 
	   } 
	 & & \\ 
 

    \theutterance \stepcounter{utterance}  

    & & & \multicolumn{2}{p{0.3\linewidth}}{\cellcolor[rgb]{0.95,0.95,0.95}{%
	\makecell[{{p{\linewidth}}}]{% 
	  \tt {\tiny [GM$|$GM]}  
	 Penalità: (4,5) non è vuoto, ma contiene '┤'.\\ \tt \\ \tt You have collectively accumulated 5 of 16 penalties.\\ \tt Riprova! 
	  } 
	   } 
	   } 
	 & & \\ 
 

    \theutterance \stepcounter{utterance}  

    & & \multicolumn{4}{p{0.6\linewidth}}{\cellcolor[rgb]{0.9,0.9,0.9}{%
	\makecell[{{p{\linewidth}}}]{% 
	  \tt {\tiny [A$\langle$GM]}  
	 Penalità: (4,5) non è vuoto, ma contiene '┤'.\\ \tt \\ \tt You have collectively accumulated 5 of 16 penalties.\\ \tt Riprova! 
	  } 
	   } 
	   } 
	 & \\ 
 

    \theutterance \stepcounter{utterance}  

    & \multicolumn{4}{p{0.6\linewidth}}{\cellcolor[rgb]{0.8,1,0.9}{%
	\makecell[{{p{\linewidth}}}]{% 
	  \tt {\tiny [A$\rangle$GM]}  
	 COMUNICA: Ho celle bloccate su x=4 alle righe y1 (D), y4 (┌), y5 (┤). Verifico visivamente che y6 e y7 su x=4 sembrano libere sul mio lato. Procedo mettendo M nella prima libera partendo da y6: tenterò M@(4,6). Confermo appena piazzato, poi proseguo con O sotto M se possibile. 
	  } 
	   } 
	   } 
	 & & \\ 
 

    \theutterance \stepcounter{utterance}  

    & & & \multicolumn{2}{p{0.3\linewidth}}{\cellcolor[rgb]{0.95,0.95,0.95}{%
	\makecell[{{p{\linewidth}}}]{% 
	  \tt {\tiny [GM$|$GM]}  
	 COMUNICA: Ho celle bloccate su x=4 alle righe y1 (D), y4 (┌), y5 (┤). Verifico visivamente che y6 e y7 su x=4 sembrano libere sul mio lato. Procedo mettendo M nella prima libera partendo da y6: tenterò M@(4,6). Confermo appena piazzato, poi proseguo con O sotto M se possibile. 
	  } 
	   } 
	   } 
	 & & \\ 
 

    \theutterance \stepcounter{utterance}  

    & & \multicolumn{4}{p{0.6\linewidth}}{\cellcolor[rgb]{0.9,0.9,0.9}{%
	\makecell[{{p{\linewidth}}}]{% 
	  \tt {\tiny [GM$\rangle$B]}  
	 Il tuo messaggio è stato inoltrato all'altro giocatore.\\ \tt \\ \tt \\ \tt You are currently playing round 13 of maximum 28.\\ \tt You have collectively accumulated 5 of 16 penalties.\\ \tt L'altro giocatore ha inviato questo messaggio:\\ \tt \\ \tt \\ \tt "Ho celle bloccate su x=4 alle righe y1 (D), y4 (┌), y5 (┤). Verifico visivamente che y6 e y7 su x=4 sembrano libere sul mio lato. Procedo mettendo M nella prima libera partendo da y6: tenterò M@(4,6). Confermo appena piazzato, poi proseguo con O sotto M se possibile."\\ \tt \\ \tt \\ \tt Qual è il tuo prossimo comando? 
	  } 
	   } 
	   } 
	 & \\ 
 

    \theutterance \stepcounter{utterance}  

    & & & \multicolumn{4}{p{0.6\linewidth}}{\cellcolor[rgb]{1,0.85,0.72}{%
	\makecell[{{p{\linewidth}}}]{% 
	  \tt {\tiny [GM$\langle$B]}  
	 SPOSTA: P, (4,4) 
	  } 
	   } 
	   } 
	  \\ 
 

    \theutterance \stepcounter{utterance}  

    & & & \multicolumn{2}{p{0.3\linewidth}}{\cellcolor[rgb]{0.95,0.95,0.95}{%
	\makecell[{{p{\linewidth}}}]{% 
	  \tt {\tiny [GM$|$GM]}  
	 SPOSTA: P, (4,4) 
	  } 
	   } 
	   } 
	 & & \\ 
 

    \theutterance \stepcounter{utterance}  

    & & & \multicolumn{2}{p{0.3\linewidth}}{\cellcolor[rgb]{0.95,0.95,0.95}{%
	\makecell[{{p{\linewidth}}}]{% 
	  \tt {\tiny [GM$|$GM]}  
	 Penalità: (4,4) non è vuoto, ma contiene '┌'.\\ \tt \\ \tt You have collectively accumulated 6 of 16 penalties.\\ \tt Riprova! 
	  } 
	   } 
	   } 
	 & & \\ 
 

    \theutterance \stepcounter{utterance}  

    & & \multicolumn{4}{p{0.6\linewidth}}{\cellcolor[rgb]{0.9,0.9,0.9}{%
	\makecell[{{p{\linewidth}}}]{% 
	  \tt {\tiny [GM$\rangle$B]}  
	 Penalità: (4,4) non è vuoto, ma contiene '┌'.\\ \tt \\ \tt You have collectively accumulated 6 of 16 penalties.\\ \tt Riprova! 
	  } 
	   } 
	   } 
	 & \\ 
 

    \theutterance \stepcounter{utterance}  

    & & & \multicolumn{4}{p{0.6\linewidth}}{\cellcolor[rgb]{1,0.85,0.72}{%
	\makecell[{{p{\linewidth}}}]{% 
	  \tt {\tiny [GM$\langle$B]}  
	 SPOSTA: P, (4,6) 
	  } 
	   } 
	   } 
	  \\ 
 

    \theutterance \stepcounter{utterance}  

    & & & \multicolumn{2}{p{0.3\linewidth}}{\cellcolor[rgb]{0.95,0.95,0.95}{%
	\makecell[{{p{\linewidth}}}]{% 
	  \tt {\tiny [GM$|$GM]}  
	 SPOSTA: P, (4,6) 
	  } 
	   } 
	   } 
	 & & \\ 
 

    \theutterance \stepcounter{utterance}  

    & & & \multicolumn{2}{p{0.3\linewidth}}{\cellcolor[rgb]{0.95,0.95,0.95}{%
	\makecell[{{p{\linewidth}}}]{% 
	  \tt {\tiny [GM$|$GM]}  
	 Penalità: (4,6) non è vuoto, ma contiene '├'.\\ \tt \\ \tt You have collectively accumulated 7 of 16 penalties.\\ \tt Riprova! 
	  } 
	   } 
	   } 
	 & & \\ 
 

    \theutterance \stepcounter{utterance}  

    & & \multicolumn{4}{p{0.6\linewidth}}{\cellcolor[rgb]{0.9,0.9,0.9}{%
	\makecell[{{p{\linewidth}}}]{% 
	  \tt {\tiny [GM$\rangle$B]}  
	 Penalità: (4,6) non è vuoto, ma contiene '├'.\\ \tt \\ \tt You have collectively accumulated 7 of 16 penalties.\\ \tt Riprova! 
	  } 
	   } 
	   } 
	 & \\ 
 

    \theutterance \stepcounter{utterance}  

    & & & \multicolumn{4}{p{0.6\linewidth}}{\cellcolor[rgb]{1,0.85,0.72}{%
	\makecell[{{p{\linewidth}}}]{% 
	  \tt {\tiny [GM$\langle$B]}  
	 SPOSTA: P, (4,7) 
	  } 
	   } 
	   } 
	  \\ 
 

    \theutterance \stepcounter{utterance}  

    & & & \multicolumn{2}{p{0.3\linewidth}}{\cellcolor[rgb]{0.95,0.95,0.95}{%
	\makecell[{{p{\linewidth}}}]{% 
	  \tt {\tiny [GM$|$GM]}  
	 SPOSTA: P, (4,7) 
	  } 
	   } 
	   } 
	 & & \\ 
 

    \theutterance \stepcounter{utterance}  

    & & & \multicolumn{2}{p{0.3\linewidth}}{\cellcolor[rgb]{0.95,0.95,0.95}{%
	\makecell[{{p{\linewidth}}}]{% 
	  \tt {\tiny [GM$|$GM]}  
	 Penalità: (4,7) non è vuoto, ma contiene '│'.\\ \tt \\ \tt You have collectively accumulated 8 of 16 penalties.\\ \tt Riprova! 
	  } 
	   } 
	   } 
	 & & \\ 
 

    \theutterance \stepcounter{utterance}  

    & & \multicolumn{4}{p{0.6\linewidth}}{\cellcolor[rgb]{0.9,0.9,0.9}{%
	\makecell[{{p{\linewidth}}}]{% 
	  \tt {\tiny [GM$\rangle$B]}  
	 Penalità: (4,7) non è vuoto, ma contiene '│'.\\ \tt \\ \tt You have collectively accumulated 8 of 16 penalties.\\ \tt Riprova! 
	  } 
	   } 
	   } 
	 & \\ 
 

    \theutterance \stepcounter{utterance}  

    & & & \multicolumn{4}{p{0.6\linewidth}}{\cellcolor[rgb]{1,0.85,0.72}{%
	\makecell[{{p{\linewidth}}}]{% 
	  \tt {\tiny [GM$\langle$B]}  
	 COMUNICA: Ho commesso errori tentando celle bloccate sulla mia x=4. Situazione corretta sul mio lato: su x=4 sono occupate da elementi di cornice/linee le righe y4 (┼ nell’intersezione), y6 (├) e y7 (│); disponibili al momento: y1, y2 con D, y3 con O, y5 è libero. Continuerò posizionando i pezzi solo in celle effettivamente libere. Aggiornami quando hai piazzato M (prevedi M@(4,6) sul tuo lato); dopo la tua conferma, io metterò P in una riga libera sotto O, compatibilmente con gli ostacoli sul mio lato. 
	  } 
	   } 
	   } 
	  \\ 
 

    \theutterance \stepcounter{utterance}  

    & & & \multicolumn{2}{p{0.3\linewidth}}{\cellcolor[rgb]{0.95,0.95,0.95}{%
	\makecell[{{p{\linewidth}}}]{% 
	  \tt {\tiny [GM$|$GM]}  
	 COMUNICA: Ho commesso errori tentando celle bloccate sulla mia x=4. Situazione corretta sul mio lato: su x=4 sono occupate da elementi di cornice/linee le righe y4 (┼ nell’intersezione), y6 (├) e y7 (│); disponibili al momento: y1, y2 con D, y3 con O, y5 è libero. Continuerò posizionando i pezzi solo in celle effettivamente libere. Aggiornami quando hai piazzato M (prevedi M@(4,6) sul tuo lato); dopo la tua conferma, io metterò P in una riga libera sotto O, compatibilmente con gli ostacoli sul mio lato. 
	  } 
	   } 
	   } 
	 & & \\ 
 

    \theutterance \stepcounter{utterance}  

    & & \multicolumn{4}{p{0.6\linewidth}}{\cellcolor[rgb]{0.9,0.9,0.9}{%
	\makecell[{{p{\linewidth}}}]{% 
	  \tt {\tiny [A$\langle$GM]}  
	 Il tuo messaggio è stato inoltrato all'altro giocatore.\\ \tt \\ \tt \\ \tt You are currently playing round 13 of maximum 28.\\ \tt You have collectively accumulated 8 of 16 penalties.\\ \tt L'altro giocatore ha inviato questo messaggio:\\ \tt \\ \tt \\ \tt "Ho commesso errori tentando celle bloccate sulla mia x=4. Situazione corretta sul mio lato: su x=4 sono occupate da elementi di cornice/linee le righe y4 (┼ nell’intersezione), y6 (├) e y7 (│); disponibili al momento: y1, y2 con D, y3 con O, y5 è libero. Continuerò posizionando i pezzi solo in celle effettivamente libere. Aggiornami quando hai piazzato M (prevedi M@(4,6) sul tuo lato); dopo la tua conferma, io metterò P in una riga libera sotto O, compatibilmente con gli ostacoli sul mio lato."\\ \tt \\ \tt \\ \tt Qual è il tuo prossimo comando? 
	  } 
	   } 
	   } 
	 & \\ 
 

    \theutterance \stepcounter{utterance}  

    & \multicolumn{4}{p{0.6\linewidth}}{\cellcolor[rgb]{0.8,1,0.9}{%
	\makecell[{{p{\linewidth}}}]{% 
	  \tt {\tiny [A$\rangle$GM]}  
	 SPOSTA: M, (4, 6) 
	  } 
	   } 
	   } 
	 & & \\ 
 

    \theutterance \stepcounter{utterance}  

    & & & \multicolumn{2}{p{0.3\linewidth}}{\cellcolor[rgb]{0.95,0.95,0.95}{%
	\makecell[{{p{\linewidth}}}]{% 
	  \tt {\tiny [GM$|$GM]}  
	 SPOSTA: M, (4, 6) 
	  } 
	   } 
	   } 
	 & & \\ 
 

    \theutterance \stepcounter{utterance}  

    & & & \multicolumn{2}{p{0.3\linewidth}}{\cellcolor[rgb]{0.95,0.95,0.95}{%
	\makecell[{{p{\linewidth}}}]{% 
	  \tt {\tiny [GM$|$GM]}  
	 Penalità: (4,6) non è vuoto, ma contiene '├'.\\ \tt \\ \tt You have collectively accumulated 9 of 16 penalties.\\ \tt Riprova! 
	  } 
	   } 
	   } 
	 & & \\ 
 

    \theutterance \stepcounter{utterance}  

    & & \multicolumn{4}{p{0.6\linewidth}}{\cellcolor[rgb]{0.9,0.9,0.9}{%
	\makecell[{{p{\linewidth}}}]{% 
	  \tt {\tiny [A$\langle$GM]}  
	 Penalità: (4,6) non è vuoto, ma contiene '├'.\\ \tt \\ \tt You have collectively accumulated 9 of 16 penalties.\\ \tt Riprova! 
	  } 
	   } 
	   } 
	 & \\ 
 

    \theutterance \stepcounter{utterance}  

    & \multicolumn{4}{p{0.6\linewidth}}{\cellcolor[rgb]{0.8,1,0.9}{%
	\makecell[{{p{\linewidth}}}]{% 
	  \tt {\tiny [A$\rangle$GM]}  
	 COMUNICA: Rivedo il mio lato: su x=4 vedo bloccate y1 (D), y4 (┌), y5 (┤); libere dovrebbero essere y2 (A), y3 (D a destra? occupata), y6 (├ bloccata), y7 (◌ libero). Quindi la prima riga libera utile per M è y7. Procedo con M@(4,7). Dopo confermo e valuto O sotto, ma potrei doverlo mettere in un’altra colonna se x=4 è esaurita. Confermi che va bene M in fondo? 
	  } 
	   } 
	   } 
	 & & \\ 
 

    \theutterance \stepcounter{utterance}  

    & & & \multicolumn{2}{p{0.3\linewidth}}{\cellcolor[rgb]{0.95,0.95,0.95}{%
	\makecell[{{p{\linewidth}}}]{% 
	  \tt {\tiny [GM$|$GM]}  
	 COMUNICA: Rivedo il mio lato: su x=4 vedo bloccate y1 (D), y4 (┌), y5 (┤); libere dovrebbero essere y2 (A), y3 (D a destra? occupata), y6 (├ bloccata), y7 (◌ libero). Quindi la prima riga libera utile per M è y7. Procedo con M@(4,7). Dopo confermo e valuto O sotto, ma potrei doverlo mettere in un’altra colonna se x=4 è esaurita. Confermi che va bene M in fondo? 
	  } 
	   } 
	   } 
	 & & \\ 
 

    \theutterance \stepcounter{utterance}  

    & & \multicolumn{4}{p{0.6\linewidth}}{\cellcolor[rgb]{0.9,0.9,0.9}{%
	\makecell[{{p{\linewidth}}}]{% 
	  \tt {\tiny [GM$\rangle$B]}  
	 Il tuo messaggio è stato inoltrato all'altro giocatore.\\ \tt \\ \tt \\ \tt You are currently playing round 14 of maximum 28.\\ \tt You have collectively accumulated 9 of 16 penalties.\\ \tt L'altro giocatore ha inviato questo messaggio:\\ \tt \\ \tt \\ \tt "Rivedo il mio lato: su x=4 vedo bloccate y1 (D), y4 (┌), y5 (┤); libere dovrebbero essere y2 (A), y3 (D a destra? occupata), y6 (├ bloccata), y7 (◌ libero). Quindi la prima riga libera utile per M è y7. Procedo con M@(4,7). Dopo confermo e valuto O sotto, ma potrei doverlo mettere in un’altra colonna se x=4 è esaurita. Confermi che va bene M in fondo?"\\ \tt \\ \tt \\ \tt Qual è il tuo prossimo comando? 
	  } 
	   } 
	   } 
	 & \\ 
 

    \theutterance \stepcounter{utterance}  

    & & & \multicolumn{4}{p{0.6\linewidth}}{\cellcolor[rgb]{1,0.85,0.72}{%
	\makecell[{{p{\linewidth}}}]{% 
	  \tt {\tiny [GM$\langle$B]}  
	 COMUNICA: Confermo, va bene M in fondo su (4,7) sul tuo lato. Sul mio lato, x=4 disponibili reali: y1 libero, y2=D, y3=O, y4 bloccato, y5 libero, y6 bloccato, y7 bloccato. Dopo la tua conferma di M@(4,7), io posizionerò P su (4,5) che è la prima riga libera sotto O, poi valuterò S e T in colonne alternative se necessario. Conferma quando M è piazzato. 
	  } 
	   } 
	   } 
	  \\ 
 

    \theutterance \stepcounter{utterance}  

    & & & \multicolumn{2}{p{0.3\linewidth}}{\cellcolor[rgb]{0.95,0.95,0.95}{%
	\makecell[{{p{\linewidth}}}]{% 
	  \tt {\tiny [GM$|$GM]}  
	 COMUNICA: Confermo, va bene M in fondo su (4,7) sul tuo lato. Sul mio lato, x=4 disponibili reali: y1 libero, y2=D, y3=O, y4 bloccato, y5 libero, y6 bloccato, y7 bloccato. Dopo la tua conferma di M@(4,7), io posizionerò P su (4,5) che è la prima riga libera sotto O, poi valuterò S e T in colonne alternative se necessario. Conferma quando M è piazzato. 
	  } 
	   } 
	   } 
	 & & \\ 
 

    \theutterance \stepcounter{utterance}  

    & & \multicolumn{4}{p{0.6\linewidth}}{\cellcolor[rgb]{0.9,0.9,0.9}{%
	\makecell[{{p{\linewidth}}}]{% 
	  \tt {\tiny [A$\langle$GM]}  
	 Il tuo messaggio è stato inoltrato all'altro giocatore.\\ \tt \\ \tt \\ \tt You are currently playing round 14 of maximum 28.\\ \tt You have collectively accumulated 9 of 16 penalties.\\ \tt L'altro giocatore ha inviato questo messaggio:\\ \tt \\ \tt \\ \tt "Confermo, va bene M in fondo su (4,7) sul tuo lato. Sul mio lato, x=4 disponibili reali: y1 libero, y2=D, y3=O, y4 bloccato, y5 libero, y6 bloccato, y7 bloccato. Dopo la tua conferma di M@(4,7), io posizionerò P su (4,5) che è la prima riga libera sotto O, poi valuterò S e T in colonne alternative se necessario. Conferma quando M è piazzato."\\ \tt \\ \tt \\ \tt Qual è il tuo prossimo comando? 
	  } 
	   } 
	   } 
	 & \\ 
 

    \theutterance \stepcounter{utterance}  

    & \multicolumn{4}{p{0.6\linewidth}}{\cellcolor[rgb]{0.8,1,0.9}{%
	\makecell[{{p{\linewidth}}}]{% 
	  \tt {\tiny [A$\rangle$GM]}  
	 SPOSTA: M, (4, 7) 
	  } 
	   } 
	   } 
	 & & \\ 
 

    \theutterance \stepcounter{utterance}  

    & & & \multicolumn{2}{p{0.3\linewidth}}{\cellcolor[rgb]{0.95,0.95,0.95}{%
	\makecell[{{p{\linewidth}}}]{% 
	  \tt {\tiny [GM$|$GM]}  
	 SPOSTA: M, (4, 7) 
	  } 
	   } 
	   } 
	 & & \\ 
 

    \theutterance \stepcounter{utterance}  

    & & & \multicolumn{2}{p{0.3\linewidth}}{\cellcolor[rgb]{0.95,0.95,0.95}{%
	\makecell[{{p{\linewidth}}}]{% 
	  \tt {\tiny [GM$|$GM]}  
	 Penalità: (4,7) non è vuoto, ma contiene '│'.\\ \tt \\ \tt You have collectively accumulated 10 of 16 penalties.\\ \tt Riprova! 
	  } 
	   } 
	   } 
	 & & \\ 
 

    \theutterance \stepcounter{utterance}  

    & & \multicolumn{4}{p{0.6\linewidth}}{\cellcolor[rgb]{0.9,0.9,0.9}{%
	\makecell[{{p{\linewidth}}}]{% 
	  \tt {\tiny [A$\langle$GM]}  
	 Penalità: (4,7) non è vuoto, ma contiene '│'.\\ \tt \\ \tt You have collectively accumulated 10 of 16 penalties.\\ \tt Riprova! 
	  } 
	   } 
	   } 
	 & \\ 
 

    \theutterance \stepcounter{utterance}  

    & \multicolumn{4}{p{0.6\linewidth}}{\cellcolor[rgb]{0.8,1,0.9}{%
	\makecell[{{p{\linewidth}}}]{% 
	  \tt {\tiny [A$\rangle$GM]}  
	 COMUNICA: Ho molte celle bloccate sulla mia colonna x=4: y1 occupata da D (iniziale), y4 da ┌, y5 da ┤, y6 da ├, y7 da │. Le uniche posizioni libere che vedo realmente sulla mia x=4 sono y2 (A già piazzato) e y3 (D già piazzato). Quindi non posso mettere altri pezzi su x=4. Propongo di spostare il nostro obiettivo su una colonna effettivamente libera su entrambi i lati: puoi dirmi quali colonne sul tuo lato hanno almeno 4 celle libere continue dall’alto verso il basso? Sul mio lato, x=2 sembra avere molte celle libere (tranne ostacoli di cornice in alcune righe), e x=5 è piena di linee; x=3 ha ostacoli; x=6 e x=7 hanno meno blocchi. Se per te x=2 è praticabile, posso riordinare lì mettendo dall’alto A,D,M,O,P,S,T nelle prime righe libere. Conferma la colonna migliore comune e ti seguo. 
	  } 
	   } 
	   } 
	 & & \\ 
 

    \theutterance \stepcounter{utterance}  

    & & & \multicolumn{2}{p{0.3\linewidth}}{\cellcolor[rgb]{0.95,0.95,0.95}{%
	\makecell[{{p{\linewidth}}}]{% 
	  \tt {\tiny [GM$|$GM]}  
	 COMUNICA: Ho molte celle bloccate sulla mia colonna x=4: y1 occupata da D (iniziale), y4 da ┌, y5 da ┤, y6 da ├, y7 da │. Le uniche posizioni libere che vedo realmente sulla mia x=4 sono y2 (A già piazzato) e y3 (D già piazzato). Quindi non posso mettere altri pezzi su x=4. Propongo di spostare il nostro obiettivo su una colonna effettivamente libera su entrambi i lati: puoi dirmi quali colonne sul tuo lato hanno almeno 4 celle libere continue dall’alto verso il basso? Sul mio lato, x=2 sembra avere molte celle libere (tranne ostacoli di cornice in alcune righe), e x=5 è piena di linee; x=3 ha ostacoli; x=6 e x=7 hanno meno blocchi. Se per te x=2 è praticabile, posso riordinare lì mettendo dall’alto A,D,M,O,P,S,T nelle prime righe libere. Conferma la colonna migliore comune e ti seguo. 
	  } 
	   } 
	   } 
	 & & \\ 
 

    \theutterance \stepcounter{utterance}  

    & & \multicolumn{4}{p{0.6\linewidth}}{\cellcolor[rgb]{0.9,0.9,0.9}{%
	\makecell[{{p{\linewidth}}}]{% 
	  \tt {\tiny [GM$\rangle$B]}  
	 Il tuo messaggio è stato inoltrato all'altro giocatore.\\ \tt \\ \tt \\ \tt You are currently playing round 15 of maximum 28.\\ \tt You have collectively accumulated 10 of 16 penalties.\\ \tt L'altro giocatore ha inviato questo messaggio:\\ \tt \\ \tt \\ \tt "Ho molte celle bloccate sulla mia colonna x=4: y1 occupata da D (iniziale), y4 da ┌, y5 da ┤, y6 da ├, y7 da │. Le uniche posizioni libere che vedo realmente sulla mia x=4 sono y2 (A già piazzato) e y3 (D già piazzato). Quindi non posso mettere altri pezzi su x=4. Propongo di spostare il nostro obiettivo su una colonna effettivamente libera su entrambi i lati: puoi dirmi quali colonne sul tuo lato hanno almeno 4 celle libere continue dall’alto verso il basso? Sul mio lato, x=2 sembra avere molte celle libere (tranne ostacoli di cornice in alcune righe), e x=5 è piena di linee; x=3 ha ostacoli; x=6 e x=7 hanno meno blocchi. Se per te x=2 è praticabile, posso riordinare lì mettendo dall’alto A,D,M,O,P,S,T nelle prime righe libere. Conferma la colonna migliore comune e ti seguo."\\ \tt \\ \tt \\ \tt Qual è il tuo prossimo comando? 
	  } 
	   } 
	   } 
	 & \\ 
 

    \theutterance \stepcounter{utterance}  

    & & & \multicolumn{4}{p{0.6\linewidth}}{\cellcolor[rgb]{1,0.85,0.72}{%
	\makecell[{{p{\linewidth}}}]{% 
	  \tt {\tiny [GM$\langle$B]}  
	 COMUNICA: Sul mio lato, panoramica colonne libere (celle realmente vuote ◌):\\ \tt - x=1: y1 libero, y2 linea, y3 A, y4 linea, y5 libero, y6 linea, y7 P → non adatta.\\ \tt - x=2: y1 S, y2 linea, y3 libero, y4 linea, y5 linea, y6 linea, y7 libero → solo 2 libere non contigue.\\ \tt - x=3: y1 libero, y2 D a destra? (libero), y3 libero, y4 linea, y5 linea, y6 linea, y7 linea → prime 3 libere dall’alto.\\ \tt - x=4: y1 libero, y2 D, y3 O, y4/6/7 bloccate, y5 libero → non adatta.\\ \tt - x=5: molte linee/ostacoli.\\ \tt - x=6: misto linee.\\ \tt - x=7: y1 libero, y2 linea, y3 libero, y4 linea, y5 libero, y6 libero, y7 libero → ha almeno 4 celle libere ma non tutte contigue in alto.\\ \tt \\ \tt Proposta: usiamo la colonna x=7 riempiendo dall’alto le prime celle libere in ordine A,D,M,O,P,S,T (saltando righe bloccate). Confermi x=7 come colonna comune con “riempi prime libere dall’alto”? Se sì, inizierò muovendo A verso la prima libera su x=7. 
	  } 
	   } 
	   } 
	  \\ 
 

    \theutterance \stepcounter{utterance}  

    & & & \multicolumn{2}{p{0.3\linewidth}}{\cellcolor[rgb]{0.95,0.95,0.95}{%
	\makecell[{{p{\linewidth}}}]{% 
	  \tt {\tiny [GM$|$GM]}  
	 COMUNICA: Sul mio lato, panoramica colonne libere (celle realmente vuote ◌):\\ \tt - x=1: y1 libero, y2 linea, y3 A, y4 linea, y5 libero, y6 linea, y7 P → non adatta.\\ \tt - x=2: y1 S, y2 linea, y3 libero, y4 linea, y5 linea, y6 linea, y7 libero → solo 2 libere non contigue.\\ \tt - x=3: y1 libero, y2 D a destra? (libero), y3 libero, y4 linea, y5 linea, y6 linea, y7 linea → prime 3 libere dall’alto.\\ \tt - x=4: y1 libero, y2 D, y3 O, y4/6/7 bloccate, y5 libero → non adatta.\\ \tt - x=5: molte linee/ostacoli.\\ \tt - x=6: misto linee.\\ \tt - x=7: y1 libero, y2 linea, y3 libero, y4 linea, y5 libero, y6 libero, y7 libero → ha almeno 4 celle libere ma non tutte contigue in alto.\\ \tt \\ \tt Proposta: usiamo la colonna x=7 riempiendo dall’alto le prime celle libere in ordine A,D,M,O,P,S,T (saltando righe bloccate). Confermi x=7 come colonna comune con “riempi prime libere dall’alto”? Se sì, inizierò muovendo A verso la prima libera su x=7. 
	  } 
	   } 
	   } 
	 & & \\ 
 

    \theutterance \stepcounter{utterance}  

    & & & \multicolumn{2}{p{0.3\linewidth}}{\cellcolor[rgb]{0.95,0.95,0.95}{%
	\makecell[{{p{\linewidth}}}]{% 
	  \tt {\tiny [GM$|$GM]}  
	 Invalid format: tail is not empty: '\\ \tt - x=1: y1 libero, y2 linea, y3 A, y4 linea, y5 libero, y6 linea, y7 P → non adatta.\\ \tt - x=2: y1 S, y2 linea, y3 libero, y4 linea, y5 linea, y6 linea, y7 libero → solo 2 libere non contigue.\\ \tt - x=3: y1 libero, y2 D a destra? (libero), y3 libero, y4 linea, y5 linea, y6 linea, y7 linea → prime 3 libere dall’alto.\\ \tt - x=4: y1 libero, y2 D, y3 O, y4/6/7 bloccate, y5 libero → non adatta.\\ \tt - x=5: molte linee/ostacoli.\\ \tt - x=6: misto linee.\\ \tt - x=7: y1 libero, y2 linea, y3 libero, y4 linea, y5 libero, y6 libero, y7 libero → ha almeno 4 celle libere ma non tutte contigue in alto.\\ \tt \\ \tt Proposta: usiamo la colonna x=7 riempiendo dall’alto le prime celle libere in ordine A,D,M,O,P,S,T (saltando righe bloccate). Confermi x=7 come colonna comune con “riempi prime libere dall’alto”? Se sì, inizierò muovendo A verso la prima libera su x=7. 
	  } 
	   } 
	   } 
	 & & \\ 
 

    \theutterance \stepcounter{utterance}  

    & & \multicolumn{4}{p{0.6\linewidth}}{\cellcolor[rgb]{0.9,0.9,0.9}{%
	\makecell[{{p{\linewidth}}}]{% 
	  \tt {\tiny [GM$\rangle$B]}  
	 Penalità: Il tuo messaggio non deve contenere nulla dopo il comando!\\ \tt \\ \tt Assicurati che la tua risposta contenga solo `COMUNICA: <MESSAGE>` o `SPOSTA: <OBJECT>, (<X>, <Y>)`, e nient'altro!\\ \tt \\ \tt You have collectively accumulated 11 of 16 penalties.\\ \tt Riprova! 
	  } 
	   } 
	   } 
	 & \\ 
 

    \theutterance \stepcounter{utterance}  

    & & & \multicolumn{4}{p{0.6\linewidth}}{\cellcolor[rgb]{1,0.85,0.72}{%
	\makecell[{{p{\linewidth}}}]{% 
	  \tt {\tiny [GM$\langle$B]}  
	 COMUNICA: Proposta nuova colonna bersaglio: x=7, riempiendo dall’alto le prime celle libere in ordine A,D,M,O,P,S,T (saltando righe bloccate). Confermi x=7? Se sì, muovo A nella prima libera su x=7. 
	  } 
	   } 
	   } 
	  \\ 
 

    \theutterance \stepcounter{utterance}  

    & & & \multicolumn{2}{p{0.3\linewidth}}{\cellcolor[rgb]{0.95,0.95,0.95}{%
	\makecell[{{p{\linewidth}}}]{% 
	  \tt {\tiny [GM$|$GM]}  
	 COMUNICA: Proposta nuova colonna bersaglio: x=7, riempiendo dall’alto le prime celle libere in ordine A,D,M,O,P,S,T (saltando righe bloccate). Confermi x=7? Se sì, muovo A nella prima libera su x=7. 
	  } 
	   } 
	   } 
	 & & \\ 
 

    \theutterance \stepcounter{utterance}  

    & & \multicolumn{4}{p{0.6\linewidth}}{\cellcolor[rgb]{0.9,0.9,0.9}{%
	\makecell[{{p{\linewidth}}}]{% 
	  \tt {\tiny [A$\langle$GM]}  
	 Il tuo messaggio è stato inoltrato all'altro giocatore.\\ \tt \\ \tt \\ \tt You are currently playing round 15 of maximum 28.\\ \tt You have collectively accumulated 11 of 16 penalties.\\ \tt L'altro giocatore ha inviato questo messaggio:\\ \tt \\ \tt \\ \tt "Proposta nuova colonna bersaglio: x=7, riempiendo dall’alto le prime celle libere in ordine A,D,M,O,P,S,T (saltando righe bloccate). Confermi x=7? Se sì, muovo A nella prima libera su x=7."\\ \tt \\ \tt \\ \tt Qual è il tuo prossimo comando? 
	  } 
	   } 
	   } 
	 & \\ 
 

    \theutterance \stepcounter{utterance}  

    & \multicolumn{4}{p{0.6\linewidth}}{\cellcolor[rgb]{0.8,1,0.9}{%
	\makecell[{{p{\linewidth}}}]{% 
	  \tt {\tiny [A$\rangle$GM]}  
	 COMUNICA: Confermo x=7 come nuova colonna bersaglio, riempiendo dall’alto le prime celle libere in ordine A,D,M,O,P,S,T e saltando righe bloccate. Procedo a spostare A nella prima cella libera su x=7 appena possibile e ti confermo la posizione; poi seguirò con D e M. 
	  } 
	   } 
	   } 
	 & & \\ 
 

    \theutterance \stepcounter{utterance}  

    & & & \multicolumn{2}{p{0.3\linewidth}}{\cellcolor[rgb]{0.95,0.95,0.95}{%
	\makecell[{{p{\linewidth}}}]{% 
	  \tt {\tiny [GM$|$GM]}  
	 COMUNICA: Confermo x=7 come nuova colonna bersaglio, riempiendo dall’alto le prime celle libere in ordine A,D,M,O,P,S,T e saltando righe bloccate. Procedo a spostare A nella prima cella libera su x=7 appena possibile e ti confermo la posizione; poi seguirò con D e M. 
	  } 
	   } 
	   } 
	 & & \\ 
 

    \theutterance \stepcounter{utterance}  

    & & \multicolumn{4}{p{0.6\linewidth}}{\cellcolor[rgb]{0.9,0.9,0.9}{%
	\makecell[{{p{\linewidth}}}]{% 
	  \tt {\tiny [GM$\rangle$B]}  
	 Il tuo messaggio è stato inoltrato all'altro giocatore.\\ \tt \\ \tt \\ \tt You are currently playing round 16 of maximum 28.\\ \tt You have collectively accumulated 11 of 16 penalties.\\ \tt L'altro giocatore ha inviato questo messaggio:\\ \tt \\ \tt \\ \tt "Confermo x=7 come nuova colonna bersaglio, riempiendo dall’alto le prime celle libere in ordine A,D,M,O,P,S,T e saltando righe bloccate. Procedo a spostare A nella prima cella libera su x=7 appena possibile e ti confermo la posizione; poi seguirò con D e M."\\ \tt \\ \tt \\ \tt Qual è il tuo prossimo comando? 
	  } 
	   } 
	   } 
	 & \\ 
 

    \theutterance \stepcounter{utterance}  

    & & & \multicolumn{4}{p{0.6\linewidth}}{\cellcolor[rgb]{1,0.85,0.72}{%
	\makecell[{{p{\linewidth}}}]{% 
	  \tt {\tiny [GM$\langle$B]}  
	 SPOSTA: A, (7,3) 
	  } 
	   } 
	   } 
	  \\ 
 

    \theutterance \stepcounter{utterance}  

    & & & \multicolumn{2}{p{0.3\linewidth}}{\cellcolor[rgb]{0.95,0.95,0.95}{%
	\makecell[{{p{\linewidth}}}]{% 
	  \tt {\tiny [GM$|$GM]}  
	 SPOSTA: A, (7,3) 
	  } 
	   } 
	   } 
	 & & \\ 
 

    \theutterance \stepcounter{utterance}  

    & & & \multicolumn{2}{p{0.3\linewidth}}{\cellcolor[rgb]{0.95,0.95,0.95}{%
	\makecell[{{p{\linewidth}}}]{% 
	  \tt {\tiny [GM$|$GM]}  
	 Spostato 'A' in (7,3) con successo. La tua griglia aggiornata appare così:\\ \tt ```\\ \tt  1234567\\ \tt ╔══╤════╗\\ \tt ║◌S│◌◌◌◌║ 1\\ \tt ╟──┘D◌┌─╢ 2\\ \tt ║◌◌◌O◌│A║ 3\\ \tt ╟─┐◌┌─┼─╢ 4\\ \tt ║◌├─┤T│◌║ 5\\ \tt ╟─┤◌├─┤◌║ 6\\ \tt ║P│◌│M│◌║ 7\\ \tt ╚═╧═╧═╧═╝\\ \tt ```\\ \tt  
	  } 
	   } 
	   } 
	 & & \\ 
 

    \theutterance \stepcounter{utterance}  

    & & \multicolumn{4}{p{0.6\linewidth}}{\cellcolor[rgb]{0.9,0.9,0.9}{%
	\makecell[{{p{\linewidth}}}]{% 
	  \tt {\tiny [A$\langle$GM]}  
	 Il tuo messaggio è stato inoltrato all'altro giocatore.\\ \tt \\ \tt \\ \tt You are currently playing round 16 of maximum 28.\\ \tt You have collectively accumulated 11 of 16 penalties.\\ \tt L'altro giocatore ha spostato un oggetto sulla sua griglia.\\ \tt \\ \tt \\ \tt Qual è il tuo prossimo comando? 
	  } 
	   } 
	   } 
	 & \\ 
 

    \theutterance \stepcounter{utterance}  

    & \multicolumn{4}{p{0.6\linewidth}}{\cellcolor[rgb]{0.8,1,0.9}{%
	\makecell[{{p{\linewidth}}}]{% 
	  \tt {\tiny [A$\rangle$GM]}  
	 SPOSTA: A, (7, 1) 
	  } 
	   } 
	   } 
	 & & \\ 
 

    \theutterance \stepcounter{utterance}  

    & & & \multicolumn{2}{p{0.3\linewidth}}{\cellcolor[rgb]{0.95,0.95,0.95}{%
	\makecell[{{p{\linewidth}}}]{% 
	  \tt {\tiny [GM$|$GM]}  
	 SPOSTA: A, (7, 1) 
	  } 
	   } 
	   } 
	 & & \\ 
 

    \theutterance \stepcounter{utterance}  

    & & & \multicolumn{2}{p{0.3\linewidth}}{\cellcolor[rgb]{0.95,0.95,0.95}{%
	\makecell[{{p{\linewidth}}}]{% 
	  \tt {\tiny [GM$|$GM]}  
	 Spostato 'A' in (7,1) con successo. La tua griglia aggiornata appare così:\\ \tt ```\\ \tt  1234567\\ \tt ╔══╤════╗\\ \tt ║◌T│◌◌◌A║ 1\\ \tt ╟──┘◌◌┌─╢ 2\\ \tt ║S◌◌D◌│◌║ 3\\ \tt ╟─┐◌┌─┼─╢ 4\\ \tt ║◌├─┤◌│O║ 5\\ \tt ╟─┤◌├─┤◌║ 6\\ \tt ║◌│P│M│◌║ 7\\ \tt ╚═╧═╧═╧═╝\\ \tt ```\\ \tt  
	  } 
	   } 
	   } 
	 & & \\ 
 

    \theutterance \stepcounter{utterance}  

    & & \multicolumn{4}{p{0.6\linewidth}}{\cellcolor[rgb]{0.9,0.9,0.9}{%
	\makecell[{{p{\linewidth}}}]{% 
	  \tt {\tiny [GM$\rangle$B]}  
	 Spostato 'A' in (7,3) con successo. La tua griglia aggiornata appare così:\\ \tt ```\\ \tt  1234567\\ \tt ╔══╤════╗\\ \tt ║◌S│◌◌◌◌║ 1\\ \tt ╟──┘D◌┌─╢ 2\\ \tt ║◌◌◌O◌│A║ 3\\ \tt ╟─┐◌┌─┼─╢ 4\\ \tt ║◌├─┤T│◌║ 5\\ \tt ╟─┤◌├─┤◌║ 6\\ \tt ║P│◌│M│◌║ 7\\ \tt ╚═╧═╧═╧═╝\\ \tt ```\\ \tt You are currently playing round 17 of maximum 28.\\ \tt You have collectively accumulated 11 of 16 penalties.\\ \tt L'altro giocatore ha spostato un oggetto sulla sua griglia.\\ \tt \\ \tt \\ \tt Qual è il tuo prossimo comando? 
	  } 
	   } 
	   } 
	 & \\ 
 

    \theutterance \stepcounter{utterance}  

    & & & \multicolumn{4}{p{0.6\linewidth}}{\cellcolor[rgb]{1,0.85,0.72}{%
	\makecell[{{p{\linewidth}}}]{% 
	  \tt {\tiny [GM$\langle$B]}  
	 COMUNICA: Ho spostato A su x=7 alla prima libera: A@(7,3). Proseguo con D nella prossima libera dall’alto su x=7; sul mio lato le celle libere su x=7 dall’alto sono y1, y2, y4, y5, y6, y7 (y3 occupata da A). Dimmi dove hai messo A sul tuo lato e se passi a D; io metterò D nella prima libera sopra A se possibile, altrimenti nella successiva libera. 
	  } 
	   } 
	   } 
	  \\ 
 

    \theutterance \stepcounter{utterance}  

    & & & \multicolumn{2}{p{0.3\linewidth}}{\cellcolor[rgb]{0.95,0.95,0.95}{%
	\makecell[{{p{\linewidth}}}]{% 
	  \tt {\tiny [GM$|$GM]}  
	 COMUNICA: Ho spostato A su x=7 alla prima libera: A@(7,3). Proseguo con D nella prossima libera dall’alto su x=7; sul mio lato le celle libere su x=7 dall’alto sono y1, y2, y4, y5, y6, y7 (y3 occupata da A). Dimmi dove hai messo A sul tuo lato e se passi a D; io metterò D nella prima libera sopra A se possibile, altrimenti nella successiva libera. 
	  } 
	   } 
	   } 
	 & & \\ 
 

    \theutterance \stepcounter{utterance}  

    & & \multicolumn{4}{p{0.6\linewidth}}{\cellcolor[rgb]{0.9,0.9,0.9}{%
	\makecell[{{p{\linewidth}}}]{% 
	  \tt {\tiny [A$\langle$GM]}  
	 Spostato 'A' in (7,1) con successo. La tua griglia aggiornata appare così:\\ \tt ```\\ \tt  1234567\\ \tt ╔══╤════╗\\ \tt ║◌T│◌◌◌A║ 1\\ \tt ╟──┘◌◌┌─╢ 2\\ \tt ║S◌◌D◌│◌║ 3\\ \tt ╟─┐◌┌─┼─╢ 4\\ \tt ║◌├─┤◌│O║ 5\\ \tt ╟─┤◌├─┤◌║ 6\\ \tt ║◌│P│M│◌║ 7\\ \tt ╚═╧═╧═╧═╝\\ \tt ```\\ \tt You are currently playing round 17 of maximum 28.\\ \tt You have collectively accumulated 11 of 16 penalties.\\ \tt L'altro giocatore ha inviato questo messaggio:\\ \tt \\ \tt \\ \tt "Ho spostato A su x=7 alla prima libera: A@(7,3). Proseguo con D nella prossima libera dall’alto su x=7; sul mio lato le celle libere su x=7 dall’alto sono y1, y2, y4, y5, y6, y7 (y3 occupata da A). Dimmi dove hai messo A sul tuo lato e se passi a D; io metterò D nella prima libera sopra A se possibile, altrimenti nella successiva libera."\\ \tt \\ \tt \\ \tt Qual è il tuo prossimo comando? 
	  } 
	   } 
	   } 
	 & \\ 
 

    \theutterance \stepcounter{utterance}  

    & \multicolumn{4}{p{0.6\linewidth}}{\cellcolor[rgb]{0.8,1,0.9}{%
	\makecell[{{p{\linewidth}}}]{% 
	  \tt {\tiny [A$\rangle$GM]}  
	 COMUNICA: Confermo A@(7,1) sul mio lato. Passo a D: tenterò di metterlo nella prima libera dall’alto su x=7; verifico che (7,2) sia libero e, se sì, lo poso lì. Dopo confermo e passo a M nella successiva libera. 
	  } 
	   } 
	   } 
	 & & \\ 
 

    \theutterance \stepcounter{utterance}  

    & & & \multicolumn{2}{p{0.3\linewidth}}{\cellcolor[rgb]{0.95,0.95,0.95}{%
	\makecell[{{p{\linewidth}}}]{% 
	  \tt {\tiny [GM$|$GM]}  
	 COMUNICA: Confermo A@(7,1) sul mio lato. Passo a D: tenterò di metterlo nella prima libera dall’alto su x=7; verifico che (7,2) sia libero e, se sì, lo poso lì. Dopo confermo e passo a M nella successiva libera. 
	  } 
	   } 
	   } 
	 & & \\ 
 

    \theutterance \stepcounter{utterance}  

    & & \multicolumn{4}{p{0.6\linewidth}}{\cellcolor[rgb]{0.9,0.9,0.9}{%
	\makecell[{{p{\linewidth}}}]{% 
	  \tt {\tiny [GM$\rangle$B]}  
	 Il tuo messaggio è stato inoltrato all'altro giocatore.\\ \tt \\ \tt \\ \tt You are currently playing round 18 of maximum 28.\\ \tt You have collectively accumulated 11 of 16 penalties.\\ \tt L'altro giocatore ha inviato questo messaggio:\\ \tt \\ \tt \\ \tt "Confermo A@(7,1) sul mio lato. Passo a D: tenterò di metterlo nella prima libera dall’alto su x=7; verifico che (7,2) sia libero e, se sì, lo poso lì. Dopo confermo e passo a M nella successiva libera."\\ \tt \\ \tt \\ \tt Qual è il tuo prossimo comando? 
	  } 
	   } 
	   } 
	 & \\ 
 

    \theutterance \stepcounter{utterance}  

    & & & \multicolumn{4}{p{0.6\linewidth}}{\cellcolor[rgb]{1,0.85,0.72}{%
	\makecell[{{p{\linewidth}}}]{% 
	  \tt {\tiny [GM$\langle$B]}  
	 SPOSTA: D, (7,1) 
	  } 
	   } 
	   } 
	  \\ 
 

    \theutterance \stepcounter{utterance}  

    & & & \multicolumn{2}{p{0.3\linewidth}}{\cellcolor[rgb]{0.95,0.95,0.95}{%
	\makecell[{{p{\linewidth}}}]{% 
	  \tt {\tiny [GM$|$GM]}  
	 SPOSTA: D, (7,1) 
	  } 
	   } 
	   } 
	 & & \\ 
 

    \theutterance \stepcounter{utterance}  

    & & & \multicolumn{2}{p{0.3\linewidth}}{\cellcolor[rgb]{0.95,0.95,0.95}{%
	\makecell[{{p{\linewidth}}}]{% 
	  \tt {\tiny [GM$|$GM]}  
	 Spostato 'D' in (7,1) con successo. La tua griglia aggiornata appare così:\\ \tt ```\\ \tt  1234567\\ \tt ╔══╤════╗\\ \tt ║◌S│◌◌◌D║ 1\\ \tt ╟──┘◌◌┌─╢ 2\\ \tt ║◌◌◌O◌│A║ 3\\ \tt ╟─┐◌┌─┼─╢ 4\\ \tt ║◌├─┤T│◌║ 5\\ \tt ╟─┤◌├─┤◌║ 6\\ \tt ║P│◌│M│◌║ 7\\ \tt ╚═╧═╧═╧═╝\\ \tt ```\\ \tt  
	  } 
	   } 
	   } 
	 & & \\ 
 

    \theutterance \stepcounter{utterance}  

    & & \multicolumn{4}{p{0.6\linewidth}}{\cellcolor[rgb]{0.9,0.9,0.9}{%
	\makecell[{{p{\linewidth}}}]{% 
	  \tt {\tiny [A$\langle$GM]}  
	 Il tuo messaggio è stato inoltrato all'altro giocatore.\\ \tt \\ \tt \\ \tt You are currently playing round 18 of maximum 28.\\ \tt You have collectively accumulated 11 of 16 penalties.\\ \tt L'altro giocatore ha spostato un oggetto sulla sua griglia.\\ \tt \\ \tt \\ \tt Qual è il tuo prossimo comando? 
	  } 
	   } 
	   } 
	 & \\ 
 

    \theutterance \stepcounter{utterance}  

    & \multicolumn{4}{p{0.6\linewidth}}{\cellcolor[rgb]{0.8,1,0.9}{%
	\makecell[{{p{\linewidth}}}]{% 
	  \tt {\tiny [A$\rangle$GM]}  
	 SPOSTA: D, (7, 2) 
	  } 
	   } 
	   } 
	 & & \\ 
 

    \theutterance \stepcounter{utterance}  

    & & & \multicolumn{2}{p{0.3\linewidth}}{\cellcolor[rgb]{0.95,0.95,0.95}{%
	\makecell[{{p{\linewidth}}}]{% 
	  \tt {\tiny [GM$|$GM]}  
	 SPOSTA: D, (7, 2) 
	  } 
	   } 
	   } 
	 & & \\ 
 

    \theutterance \stepcounter{utterance}  

    & & & \multicolumn{2}{p{0.3\linewidth}}{\cellcolor[rgb]{0.95,0.95,0.95}{%
	\makecell[{{p{\linewidth}}}]{% 
	  \tt {\tiny [GM$|$GM]}  
	 Penalità: (7,2) non è vuoto, ma contiene '─'.\\ \tt \\ \tt You have collectively accumulated 12 of 16 penalties.\\ \tt Riprova! 
	  } 
	   } 
	   } 
	 & & \\ 
 

    \theutterance \stepcounter{utterance}  

    & & \multicolumn{4}{p{0.6\linewidth}}{\cellcolor[rgb]{0.9,0.9,0.9}{%
	\makecell[{{p{\linewidth}}}]{% 
	  \tt {\tiny [A$\langle$GM]}  
	 Penalità: (7,2) non è vuoto, ma contiene '─'.\\ \tt \\ \tt You have collectively accumulated 12 of 16 penalties.\\ \tt Riprova! 
	  } 
	   } 
	   } 
	 & \\ 
 

    \theutterance \stepcounter{utterance}  

    & \multicolumn{4}{p{0.6\linewidth}}{\cellcolor[rgb]{0.8,1,0.9}{%
	\makecell[{{p{\linewidth}}}]{% 
	  \tt {\tiny [A$\rangle$GM]}  
	 COMUNICA: Sul mio lato (7,2) è bloccato da una linea. Cerco la prima riga libera su x=7 partendo dall’alto dopo A@(7,1). Visivamente sono libere y4, y6; y3 è occupata da ostacoli sulla mia griglia. Procedo mettendo D nella prima davvero libera: tenterò D@(7,4). Se va a buon fine, poi proverò M nella successiva libera (probabilmente y6). Confermo appena piazzato. 
	  } 
	   } 
	   } 
	 & & \\ 
 

    \theutterance \stepcounter{utterance}  

    & & & \multicolumn{2}{p{0.3\linewidth}}{\cellcolor[rgb]{0.95,0.95,0.95}{%
	\makecell[{{p{\linewidth}}}]{% 
	  \tt {\tiny [GM$|$GM]}  
	 COMUNICA: Sul mio lato (7,2) è bloccato da una linea. Cerco la prima riga libera su x=7 partendo dall’alto dopo A@(7,1). Visivamente sono libere y4, y6; y3 è occupata da ostacoli sulla mia griglia. Procedo mettendo D nella prima davvero libera: tenterò D@(7,4). Se va a buon fine, poi proverò M nella successiva libera (probabilmente y6). Confermo appena piazzato. 
	  } 
	   } 
	   } 
	 & & \\ 
 

    \theutterance \stepcounter{utterance}  

    & & \multicolumn{4}{p{0.6\linewidth}}{\cellcolor[rgb]{0.9,0.9,0.9}{%
	\makecell[{{p{\linewidth}}}]{% 
	  \tt {\tiny [GM$\rangle$B]}  
	 Spostato 'D' in (7,1) con successo. La tua griglia aggiornata appare così:\\ \tt ```\\ \tt  1234567\\ \tt ╔══╤════╗\\ \tt ║◌S│◌◌◌D║ 1\\ \tt ╟──┘◌◌┌─╢ 2\\ \tt ║◌◌◌O◌│A║ 3\\ \tt ╟─┐◌┌─┼─╢ 4\\ \tt ║◌├─┤T│◌║ 5\\ \tt ╟─┤◌├─┤◌║ 6\\ \tt ║P│◌│M│◌║ 7\\ \tt ╚═╧═╧═╧═╝\\ \tt ```\\ \tt You are currently playing round 19 of maximum 28.\\ \tt You have collectively accumulated 12 of 16 penalties.\\ \tt L'altro giocatore ha inviato questo messaggio:\\ \tt \\ \tt \\ \tt "Sul mio lato (7,2) è bloccato da una linea. Cerco la prima riga libera su x=7 partendo dall’alto dopo A@(7,1). Visivamente sono libere y4, y6; y3 è occupata da ostacoli sulla mia griglia. Procedo mettendo D nella prima davvero libera: tenterò D@(7,4). Se va a buon fine, poi proverò M nella successiva libera (probabilmente y6). Confermo appena piazzato."\\ \tt \\ \tt \\ \tt Qual è il tuo prossimo comando? 
	  } 
	   } 
	   } 
	 & \\ 
 

    \theutterance \stepcounter{utterance}  

    & & & \multicolumn{4}{p{0.6\linewidth}}{\cellcolor[rgb]{1,0.85,0.72}{%
	\makecell[{{p{\linewidth}}}]{% 
	  \tt {\tiny [GM$\langle$B]}  
	 COMUNICA: Ricevuto. Sul mio lato situazione x=7 dall’alto: y1=D, y2 libero, y3=A, y4 libero, y5 libero, y6 libero, y7 libero. Procedo a mettere M nella prima libera disponibile sopra A se serve mantenere ordine dall’alto, altrimenti sotto A seguendo A,D,M,O,P,S,T. Conferma quando hai piazzato D@(7,4) e la tua prossima per M. 
	  } 
	   } 
	   } 
	  \\ 
 

    \theutterance \stepcounter{utterance}  

    & & & \multicolumn{2}{p{0.3\linewidth}}{\cellcolor[rgb]{0.95,0.95,0.95}{%
	\makecell[{{p{\linewidth}}}]{% 
	  \tt {\tiny [GM$|$GM]}  
	 COMUNICA: Ricevuto. Sul mio lato situazione x=7 dall’alto: y1=D, y2 libero, y3=A, y4 libero, y5 libero, y6 libero, y7 libero. Procedo a mettere M nella prima libera disponibile sopra A se serve mantenere ordine dall’alto, altrimenti sotto A seguendo A,D,M,O,P,S,T. Conferma quando hai piazzato D@(7,4) e la tua prossima per M. 
	  } 
	   } 
	   } 
	 & & \\ 
 

    \theutterance \stepcounter{utterance}  

    & & \multicolumn{4}{p{0.6\linewidth}}{\cellcolor[rgb]{0.9,0.9,0.9}{%
	\makecell[{{p{\linewidth}}}]{% 
	  \tt {\tiny [A$\langle$GM]}  
	 Il tuo messaggio è stato inoltrato all'altro giocatore.\\ \tt \\ \tt \\ \tt You are currently playing round 19 of maximum 28.\\ \tt You have collectively accumulated 12 of 16 penalties.\\ \tt L'altro giocatore ha inviato questo messaggio:\\ \tt \\ \tt \\ \tt "Ricevuto. Sul mio lato situazione x=7 dall’alto: y1=D, y2 libero, y3=A, y4 libero, y5 libero, y6 libero, y7 libero. Procedo a mettere M nella prima libera disponibile sopra A se serve mantenere ordine dall’alto, altrimenti sotto A seguendo A,D,M,O,P,S,T. Conferma quando hai piazzato D@(7,4) e la tua prossima per M."\\ \tt \\ \tt \\ \tt Qual è il tuo prossimo comando? 
	  } 
	   } 
	   } 
	 & \\ 
 

    \theutterance \stepcounter{utterance}  

    & \multicolumn{4}{p{0.6\linewidth}}{\cellcolor[rgb]{0.8,1,0.9}{%
	\makecell[{{p{\linewidth}}}]{% 
	  \tt {\tiny [A$\rangle$GM]}  
	 SPOSTA: D, (7, 4) 
	  } 
	   } 
	   } 
	 & & \\ 
 

    \theutterance \stepcounter{utterance}  

    & & & \multicolumn{2}{p{0.3\linewidth}}{\cellcolor[rgb]{0.95,0.95,0.95}{%
	\makecell[{{p{\linewidth}}}]{% 
	  \tt {\tiny [GM$|$GM]}  
	 SPOSTA: D, (7, 4) 
	  } 
	   } 
	   } 
	 & & \\ 
 

    \theutterance \stepcounter{utterance}  

    & & & \multicolumn{2}{p{0.3\linewidth}}{\cellcolor[rgb]{0.95,0.95,0.95}{%
	\makecell[{{p{\linewidth}}}]{% 
	  \tt {\tiny [GM$|$GM]}  
	 Penalità: (7,4) non è vuoto, ma contiene '─'.\\ \tt \\ \tt You have collectively accumulated 13 of 16 penalties.\\ \tt Riprova! 
	  } 
	   } 
	   } 
	 & & \\ 
 

    \theutterance \stepcounter{utterance}  

    & & \multicolumn{4}{p{0.6\linewidth}}{\cellcolor[rgb]{0.9,0.9,0.9}{%
	\makecell[{{p{\linewidth}}}]{% 
	  \tt {\tiny [A$\langle$GM]}  
	 Penalità: (7,4) non è vuoto, ma contiene '─'.\\ \tt \\ \tt You have collectively accumulated 13 of 16 penalties.\\ \tt Riprova! 
	  } 
	   } 
	   } 
	 & \\ 
 

    \theutterance \stepcounter{utterance}  

    & \multicolumn{4}{p{0.6\linewidth}}{\cellcolor[rgb]{0.8,1,0.9}{%
	\makecell[{{p{\linewidth}}}]{% 
	  \tt {\tiny [A$\rangle$GM]}  
	 COMUNICA: Ho verificato meglio: sulla mia colonna x=7 le righe y2 e y4 sono bloccate da linee, y5 è libera, y6 libera, y7 libera. Propongo di posizionare D nella prima davvero libera sotto A, cioè D@(7,5) sul mio lato; poi M@(7,6). Tu mantieni l’ordine A,D,M,O,P,S,T dall’alto dove possibile sulle tue celle libere. Confermo appena riesco a mettere D su (7,5). 
	  } 
	   } 
	   } 
	 & & \\ 
 

    \theutterance \stepcounter{utterance}  

    & & & \multicolumn{2}{p{0.3\linewidth}}{\cellcolor[rgb]{0.95,0.95,0.95}{%
	\makecell[{{p{\linewidth}}}]{% 
	  \tt {\tiny [GM$|$GM]}  
	 COMUNICA: Ho verificato meglio: sulla mia colonna x=7 le righe y2 e y4 sono bloccate da linee, y5 è libera, y6 libera, y7 libera. Propongo di posizionare D nella prima davvero libera sotto A, cioè D@(7,5) sul mio lato; poi M@(7,6). Tu mantieni l’ordine A,D,M,O,P,S,T dall’alto dove possibile sulle tue celle libere. Confermo appena riesco a mettere D su (7,5). 
	  } 
	   } 
	   } 
	 & & \\ 
 

    \theutterance \stepcounter{utterance}  

    & & \multicolumn{4}{p{0.6\linewidth}}{\cellcolor[rgb]{0.9,0.9,0.9}{%
	\makecell[{{p{\linewidth}}}]{% 
	  \tt {\tiny [GM$\rangle$B]}  
	 Il tuo messaggio è stato inoltrato all'altro giocatore.\\ \tt \\ \tt \\ \tt You are currently playing round 20 of maximum 28.\\ \tt You have collectively accumulated 13 of 16 penalties.\\ \tt L'altro giocatore ha inviato questo messaggio:\\ \tt \\ \tt \\ \tt "Ho verificato meglio: sulla mia colonna x=7 le righe y2 e y4 sono bloccate da linee, y5 è libera, y6 libera, y7 libera. Propongo di posizionare D nella prima davvero libera sotto A, cioè D@(7,5) sul mio lato; poi M@(7,6). Tu mantieni l’ordine A,D,M,O,P,S,T dall’alto dove possibile sulle tue celle libere. Confermo appena riesco a mettere D su (7,5)."\\ \tt \\ \tt \\ \tt Qual è il tuo prossimo comando? 
	  } 
	   } 
	   } 
	 & \\ 
 

    \theutterance \stepcounter{utterance}  

    & & & \multicolumn{4}{p{0.6\linewidth}}{\cellcolor[rgb]{1,0.85,0.72}{%
	\makecell[{{p{\linewidth}}}]{% 
	  \tt {\tiny [GM$\langle$B]}  
	 SPOSTA: M, (7,2) 
	  } 
	   } 
	   } 
	  \\ 
 

    \theutterance \stepcounter{utterance}  

    & & & \multicolumn{2}{p{0.3\linewidth}}{\cellcolor[rgb]{0.95,0.95,0.95}{%
	\makecell[{{p{\linewidth}}}]{% 
	  \tt {\tiny [GM$|$GM]}  
	 SPOSTA: M, (7,2) 
	  } 
	   } 
	   } 
	 & & \\ 
 

    \theutterance \stepcounter{utterance}  

    & & & \multicolumn{2}{p{0.3\linewidth}}{\cellcolor[rgb]{0.95,0.95,0.95}{%
	\makecell[{{p{\linewidth}}}]{% 
	  \tt {\tiny [GM$|$GM]}  
	 Penalità: (7,2) non è vuoto, ma contiene '─'.\\ \tt \\ \tt You have collectively accumulated 14 of 16 penalties.\\ \tt Riprova! 
	  } 
	   } 
	   } 
	 & & \\ 
 

    \theutterance \stepcounter{utterance}  

    & & \multicolumn{4}{p{0.6\linewidth}}{\cellcolor[rgb]{0.9,0.9,0.9}{%
	\makecell[{{p{\linewidth}}}]{% 
	  \tt {\tiny [GM$\rangle$B]}  
	 Penalità: (7,2) non è vuoto, ma contiene '─'.\\ \tt \\ \tt You have collectively accumulated 14 of 16 penalties.\\ \tt Riprova! 
	  } 
	   } 
	   } 
	 & \\ 
 

    \theutterance \stepcounter{utterance}  

    & & & \multicolumn{4}{p{0.6\linewidth}}{\cellcolor[rgb]{1,0.85,0.72}{%
	\makecell[{{p{\linewidth}}}]{% 
	  \tt {\tiny [GM$\langle$B]}  
	 SPOSTA: M, (7,4) 
	  } 
	   } 
	   } 
	  \\ 
 

    \theutterance \stepcounter{utterance}  

    & & & \multicolumn{2}{p{0.3\linewidth}}{\cellcolor[rgb]{0.95,0.95,0.95}{%
	\makecell[{{p{\linewidth}}}]{% 
	  \tt {\tiny [GM$|$GM]}  
	 SPOSTA: M, (7,4) 
	  } 
	   } 
	   } 
	 & & \\ 
 

    \theutterance \stepcounter{utterance}  

    & & & \multicolumn{2}{p{0.3\linewidth}}{\cellcolor[rgb]{0.95,0.95,0.95}{%
	\makecell[{{p{\linewidth}}}]{% 
	  \tt {\tiny [GM$|$GM]}  
	 Penalità: (7,4) non è vuoto, ma contiene '─'.\\ \tt \\ \tt You have collectively accumulated 15 of 16 penalties.\\ \tt Riprova! 
	  } 
	   } 
	   } 
	 & & \\ 
 

    \theutterance \stepcounter{utterance}  

    & & \multicolumn{4}{p{0.6\linewidth}}{\cellcolor[rgb]{0.9,0.9,0.9}{%
	\makecell[{{p{\linewidth}}}]{% 
	  \tt {\tiny [GM$\rangle$B]}  
	 Penalità: (7,4) non è vuoto, ma contiene '─'.\\ \tt \\ \tt You have collectively accumulated 15 of 16 penalties.\\ \tt Riprova! 
	  } 
	   } 
	   } 
	 & \\ 
 

    \theutterance \stepcounter{utterance}  

    & & & \multicolumn{4}{p{0.6\linewidth}}{\cellcolor[rgb]{1,0.85,0.72}{%
	\makecell[{{p{\linewidth}}}]{% 
	  \tt {\tiny [GM$\langle$B]}  
	 COMUNICA: Ho commesso errori sulle celle bloccate. Stato corretto colonna x=7 sul mio lato: y1=D, y2 bloccata (─), y3=A, y4 bloccata (─), y5 libero, y6 libero, y7 libero. Attendo tua conferma quando metti D@(7,5); poi io metterò M nella prima libera successiva, cioè (7,6). 
	  } 
	   } 
	   } 
	  \\ 
 

    \theutterance \stepcounter{utterance}  

    & & & \multicolumn{2}{p{0.3\linewidth}}{\cellcolor[rgb]{0.95,0.95,0.95}{%
	\makecell[{{p{\linewidth}}}]{% 
	  \tt {\tiny [GM$|$GM]}  
	 COMUNICA: Ho commesso errori sulle celle bloccate. Stato corretto colonna x=7 sul mio lato: y1=D, y2 bloccata (─), y3=A, y4 bloccata (─), y5 libero, y6 libero, y7 libero. Attendo tua conferma quando metti D@(7,5); poi io metterò M nella prima libera successiva, cioè (7,6). 
	  } 
	   } 
	   } 
	 & & \\ 
 

    \theutterance \stepcounter{utterance}  

    & & \multicolumn{4}{p{0.6\linewidth}}{\cellcolor[rgb]{0.9,0.9,0.9}{%
	\makecell[{{p{\linewidth}}}]{% 
	  \tt {\tiny [A$\langle$GM]}  
	 Il tuo messaggio è stato inoltrato all'altro giocatore.\\ \tt \\ \tt \\ \tt You are currently playing round 20 of maximum 28.\\ \tt You have collectively accumulated 15 of 16 penalties.\\ \tt L'altro giocatore ha inviato questo messaggio:\\ \tt \\ \tt \\ \tt "Ho commesso errori sulle celle bloccate. Stato corretto colonna x=7 sul mio lato: y1=D, y2 bloccata (─), y3=A, y4 bloccata (─), y5 libero, y6 libero, y7 libero. Attendo tua conferma quando metti D@(7,5); poi io metterò M nella prima libera successiva, cioè (7,6)."\\ \tt \\ \tt \\ \tt Qual è il tuo prossimo comando? 
	  } 
	   } 
	   } 
	 & \\ 
 

    \theutterance \stepcounter{utterance}  

    & \multicolumn{4}{p{0.6\linewidth}}{\cellcolor[rgb]{0.8,1,0.9}{%
	\makecell[{{p{\linewidth}}}]{% 
	  \tt {\tiny [A$\rangle$GM]}  
	 SPOSTA: D, (7, 5) 
	  } 
	   } 
	   } 
	 & & \\ 
 

    \theutterance \stepcounter{utterance}  

    & & & \multicolumn{2}{p{0.3\linewidth}}{\cellcolor[rgb]{0.95,0.95,0.95}{%
	\makecell[{{p{\linewidth}}}]{% 
	  \tt {\tiny [GM$|$GM]}  
	 SPOSTA: D, (7, 5) 
	  } 
	   } 
	   } 
	 & & \\ 
 

    \theutterance \stepcounter{utterance}  

    & & & \multicolumn{2}{p{0.3\linewidth}}{\cellcolor[rgb]{0.95,0.95,0.95}{%
	\makecell[{{p{\linewidth}}}]{% 
	  \tt {\tiny [GM$|$GM]}  
	 Penalità: (7,5) non è vuoto, ma contiene 'O'.\\ \tt \\ \tt You have collectively accumulated 16 of 16 penalties.\\ \tt Riprova! 
	  } 
	   } 
	   } 
	 & & \\ 
 

    \theutterance \stepcounter{utterance}  

    & & \multicolumn{4}{p{0.6\linewidth}}{\cellcolor[rgb]{0.9,0.9,0.9}{%
	\makecell[{{p{\linewidth}}}]{% 
	  \tt {\tiny [A$\langle$GM]}  
	 Penalità: (7,5) non è vuoto, ma contiene 'O'.\\ \tt \\ \tt You have collectively accumulated 16 of 16 penalties.\\ \tt Riprova! 
	  } 
	   } 
	   } 
	 & \\ 
 

    \theutterance \stepcounter{utterance}  

    & \multicolumn{4}{p{0.6\linewidth}}{\cellcolor[rgb]{0.8,1,0.9}{%
	\makecell[{{p{\linewidth}}}]{% 
	  \tt {\tiny [A$\rangle$GM]}  
	 COMUNICA: Purtroppo ho appena tentato una cella occupata: sulla mia griglia (7,5) è già occupata da O. Il nostro contatore penalità ha raggiunto il limite; temo che la partita sia compromessa. Se possiamo ancora fare un’ultima azione utile: confermo stato mio su x=7 — A@(7,1), O@(7,5), libere y6–y7 — e proverei D@(7,6) come prossima libera, poi M@(7,7). Dimmi se vuoi comunque tentare di chiudere con questo ordine residuo. 
	  } 
	   } 
	   } 
	 & & \\ 
 

    \theutterance \stepcounter{utterance}  

    & & & \multicolumn{2}{p{0.3\linewidth}}{\cellcolor[rgb]{0.95,0.95,0.95}{%
	\makecell[{{p{\linewidth}}}]{% 
	  \tt {\tiny [GM$|$GM]}  
	 COMUNICA: Purtroppo ho appena tentato una cella occupata: sulla mia griglia (7,5) è già occupata da O. Il nostro contatore penalità ha raggiunto il limite; temo che la partita sia compromessa. Se possiamo ancora fare un’ultima azione utile: confermo stato mio su x=7 — A@(7,1), O@(7,5), libere y6–y7 — e proverei D@(7,6) come prossima libera, poi M@(7,7). Dimmi se vuoi comunque tentare di chiudere con questo ordine residuo. 
	  } 
	   } 
	   } 
	 & & \\ 
 

    \theutterance \stepcounter{utterance}  

    & & \multicolumn{4}{p{0.6\linewidth}}{\cellcolor[rgb]{0.9,0.9,0.9}{%
	\makecell[{{p{\linewidth}}}]{% 
	  \tt {\tiny [GM$\rangle$B]}  
	 Il tuo messaggio è stato inoltrato all'altro giocatore.\\ \tt \\ \tt \\ \tt You are currently playing round 21 of maximum 28.\\ \tt You have collectively accumulated 16 of 16 penalties.\\ \tt L'altro giocatore ha inviato questo messaggio:\\ \tt \\ \tt \\ \tt "Purtroppo ho appena tentato una cella occupata: sulla mia griglia (7,5) è già occupata da O. Il nostro contatore penalità ha raggiunto il limite; temo che la partita sia compromessa. Se possiamo ancora fare un’ultima azione utile: confermo stato mio su x=7 — A@(7,1), O@(7,5), libere y6–y7 — e proverei D@(7,6) come prossima libera, poi M@(7,7). Dimmi se vuoi comunque tentare di chiudere con questo ordine residuo."\\ \tt \\ \tt \\ \tt Qual è il tuo prossimo comando? 
	  } 
	   } 
	   } 
	 & \\ 
 

    \theutterance \stepcounter{utterance}  

    & & & \multicolumn{4}{p{0.6\linewidth}}{\cellcolor[rgb]{1,0.85,0.72}{%
	\makecell[{{p{\linewidth}}}]{% 
	  \tt {\tiny [GM$\langle$B]}  
	 COMUNICA: Capito. Anche se abbiamo raggiunto il limite penalità, tentiamo comunque di allinearci il più possibile: sul mio lato x=7 ho D@(7,1), A@(7,3), libere y5–y7. Procedi pure con D@(7,6) e M@(7,7) sul tuo lato; io posizionerò M@(7,6) e poi O@(7,7) se restano libere, per mantenere l’ordine A,D,M,O,P,S,T dall’alto. Conferma esito. 
	  } 
	   } 
	   } 
	  \\ 
 

    \theutterance \stepcounter{utterance}  

    & & & \multicolumn{2}{p{0.3\linewidth}}{\cellcolor[rgb]{0.95,0.95,0.95}{%
	\makecell[{{p{\linewidth}}}]{% 
	  \tt {\tiny [GM$|$GM]}  
	 COMUNICA: Capito. Anche se abbiamo raggiunto il limite penalità, tentiamo comunque di allinearci il più possibile: sul mio lato x=7 ho D@(7,1), A@(7,3), libere y5–y7. Procedi pure con D@(7,6) e M@(7,7) sul tuo lato; io posizionerò M@(7,6) e poi O@(7,7) se restano libere, per mantenere l’ordine A,D,M,O,P,S,T dall’alto. Conferma esito. 
	  } 
	   } 
	   } 
	 & & \\ 
 

    \theutterance \stepcounter{utterance}  

    & & \multicolumn{4}{p{0.6\linewidth}}{\cellcolor[rgb]{0.9,0.9,0.9}{%
	\makecell[{{p{\linewidth}}}]{% 
	  \tt {\tiny [A$\langle$GM]}  
	 Il tuo messaggio è stato inoltrato all'altro giocatore.\\ \tt \\ \tt \\ \tt You are currently playing round 21 of maximum 28.\\ \tt You have collectively accumulated 16 of 16 penalties.\\ \tt L'altro giocatore ha inviato questo messaggio:\\ \tt \\ \tt \\ \tt "Capito. Anche se abbiamo raggiunto il limite penalità, tentiamo comunque di allinearci il più possibile: sul mio lato x=7 ho D@(7,1), A@(7,3), libere y5–y7. Procedi pure con D@(7,6) e M@(7,7) sul tuo lato; io posizionerò M@(7,6) e poi O@(7,7) se restano libere, per mantenere l’ordine A,D,M,O,P,S,T dall’alto. Conferma esito."\\ \tt \\ \tt \\ \tt Qual è il tuo prossimo comando? 
	  } 
	   } 
	   } 
	 & \\ 
 

    \theutterance \stepcounter{utterance}  

    & \multicolumn{4}{p{0.6\linewidth}}{\cellcolor[rgb]{0.8,1,0.9}{%
	\makecell[{{p{\linewidth}}}]{% 
	  \tt {\tiny [A$\rangle$GM]}  
	 SPOSTA: D, (7, 6) 
	  } 
	   } 
	   } 
	 & & \\ 
 

    \theutterance \stepcounter{utterance}  

    & & & \multicolumn{2}{p{0.3\linewidth}}{\cellcolor[rgb]{0.95,0.95,0.95}{%
	\makecell[{{p{\linewidth}}}]{% 
	  \tt {\tiny [GM$|$GM]}  
	 SPOSTA: D, (7, 6) 
	  } 
	   } 
	   } 
	 & & \\ 
 

    \theutterance \stepcounter{utterance}  

    & & & \multicolumn{2}{p{0.3\linewidth}}{\cellcolor[rgb]{0.95,0.95,0.95}{%
	\makecell[{{p{\linewidth}}}]{% 
	  \tt {\tiny [GM$|$GM]}  
	 Spostato 'D' in (7,6) con successo. La tua griglia aggiornata appare così:\\ \tt ```\\ \tt  1234567\\ \tt ╔══╤════╗\\ \tt ║◌T│◌◌◌A║ 1\\ \tt ╟──┘◌◌┌─╢ 2\\ \tt ║S◌◌◌◌│◌║ 3\\ \tt ╟─┐◌┌─┼─╢ 4\\ \tt ║◌├─┤◌│O║ 5\\ \tt ╟─┤◌├─┤D║ 6\\ \tt ║◌│P│M│◌║ 7\\ \tt ╚═╧═╧═╧═╝\\ \tt ```\\ \tt  
	  } 
	   } 
	   } 
	 & & \\ 
 

    \theutterance \stepcounter{utterance}  

    & & \multicolumn{4}{p{0.6\linewidth}}{\cellcolor[rgb]{0.9,0.9,0.9}{%
	\makecell[{{p{\linewidth}}}]{% 
	  \tt {\tiny [GM$\rangle$B]}  
	 Il tuo messaggio è stato inoltrato all'altro giocatore.\\ \tt \\ \tt \\ \tt You are currently playing round 22 of maximum 28.\\ \tt You have collectively accumulated 16 of 16 penalties.\\ \tt L'altro giocatore ha spostato un oggetto sulla sua griglia.\\ \tt \\ \tt \\ \tt Qual è il tuo prossimo comando? 
	  } 
	   } 
	   } 
	 & \\ 
 

    \theutterance \stepcounter{utterance}  

    & & & \multicolumn{4}{p{0.6\linewidth}}{\cellcolor[rgb]{1,0.85,0.72}{%
	\makecell[{{p{\linewidth}}}]{% 
	  \tt {\tiny [GM$\langle$B]}  
	 SPOSTA: M, (7,6) 
	  } 
	   } 
	   } 
	  \\ 
 

    \theutterance \stepcounter{utterance}  

    & & & \multicolumn{2}{p{0.3\linewidth}}{\cellcolor[rgb]{0.95,0.95,0.95}{%
	\makecell[{{p{\linewidth}}}]{% 
	  \tt {\tiny [GM$|$GM]}  
	 SPOSTA: M, (7,6) 
	  } 
	   } 
	   } 
	 & & \\ 
 

    \theutterance \stepcounter{utterance}  

    & & & \multicolumn{2}{p{0.3\linewidth}}{\cellcolor[rgb]{0.95,0.95,0.95}{%
	\makecell[{{p{\linewidth}}}]{% 
	  \tt {\tiny [GM$|$GM]}  
	 Spostato 'M' in (7,6) con successo. La tua griglia aggiornata appare così:\\ \tt ```\\ \tt  1234567\\ \tt ╔══╤════╗\\ \tt ║◌S│◌◌◌D║ 1\\ \tt ╟──┘◌◌┌─╢ 2\\ \tt ║◌◌◌O◌│A║ 3\\ \tt ╟─┐◌┌─┼─╢ 4\\ \tt ║◌├─┤T│◌║ 5\\ \tt ╟─┤◌├─┤M║ 6\\ \tt ║P│◌│◌│◌║ 7\\ \tt ╚═╧═╧═╧═╝\\ \tt ```\\ \tt  
	  } 
	   } 
	   } 
	 & & \\ 
 

    \theutterance \stepcounter{utterance}  

    & & \multicolumn{4}{p{0.6\linewidth}}{\cellcolor[rgb]{0.9,0.9,0.9}{%
	\makecell[{{p{\linewidth}}}]{% 
	  \tt {\tiny [A$\langle$GM]}  
	 Spostato 'D' in (7,6) con successo. La tua griglia aggiornata appare così:\\ \tt ```\\ \tt  1234567\\ \tt ╔══╤════╗\\ \tt ║◌T│◌◌◌A║ 1\\ \tt ╟──┘◌◌┌─╢ 2\\ \tt ║S◌◌◌◌│◌║ 3\\ \tt ╟─┐◌┌─┼─╢ 4\\ \tt ║◌├─┤◌│O║ 5\\ \tt ╟─┤◌├─┤D║ 6\\ \tt ║◌│P│M│◌║ 7\\ \tt ╚═╧═╧═╧═╝\\ \tt ```\\ \tt You are currently playing round 22 of maximum 28.\\ \tt You have collectively accumulated 16 of 16 penalties.\\ \tt L'altro giocatore ha spostato un oggetto sulla sua griglia.\\ \tt \\ \tt \\ \tt Qual è il tuo prossimo comando? 
	  } 
	   } 
	   } 
	 & \\ 
 

    \theutterance \stepcounter{utterance}  

    & \multicolumn{4}{p{0.6\linewidth}}{\cellcolor[rgb]{0.8,1,0.9}{%
	\makecell[{{p{\linewidth}}}]{% 
	  \tt {\tiny [A$\rangle$GM]}  
	 SPOSTA: M, (7, 7) 
	  } 
	   } 
	   } 
	 & & \\ 
 

    \theutterance \stepcounter{utterance}  

    & & & \multicolumn{2}{p{0.3\linewidth}}{\cellcolor[rgb]{0.95,0.95,0.95}{%
	\makecell[{{p{\linewidth}}}]{% 
	  \tt {\tiny [GM$|$GM]}  
	 SPOSTA: M, (7, 7) 
	  } 
	   } 
	   } 
	 & & \\ 
 

    \theutterance \stepcounter{utterance}  

    & & & \multicolumn{2}{p{0.3\linewidth}}{\cellcolor[rgb]{0.95,0.95,0.95}{%
	\makecell[{{p{\linewidth}}}]{% 
	  \tt {\tiny [GM$|$GM]}  
	 Spostato 'M' in (7,7) con successo. La tua griglia aggiornata appare così:\\ \tt ```\\ \tt  1234567\\ \tt ╔══╤════╗\\ \tt ║◌T│◌◌◌A║ 1\\ \tt ╟──┘◌◌┌─╢ 2\\ \tt ║S◌◌◌◌│◌║ 3\\ \tt ╟─┐◌┌─┼─╢ 4\\ \tt ║◌├─┤◌│O║ 5\\ \tt ╟─┤◌├─┤D║ 6\\ \tt ║◌│P│◌│M║ 7\\ \tt ╚═╧═╧═╧═╝\\ \tt ```\\ \tt  
	  } 
	   } 
	   } 
	 & & \\ 
 

    \theutterance \stepcounter{utterance}  

    & & \multicolumn{4}{p{0.6\linewidth}}{\cellcolor[rgb]{0.9,0.9,0.9}{%
	\makecell[{{p{\linewidth}}}]{% 
	  \tt {\tiny [GM$\rangle$B]}  
	 Spostato 'M' in (7,6) con successo. La tua griglia aggiornata appare così:\\ \tt ```\\ \tt  1234567\\ \tt ╔══╤════╗\\ \tt ║◌S│◌◌◌D║ 1\\ \tt ╟──┘◌◌┌─╢ 2\\ \tt ║◌◌◌O◌│A║ 3\\ \tt ╟─┐◌┌─┼─╢ 4\\ \tt ║◌├─┤T│◌║ 5\\ \tt ╟─┤◌├─┤M║ 6\\ \tt ║P│◌│◌│◌║ 7\\ \tt ╚═╧═╧═╧═╝\\ \tt ```\\ \tt You are currently playing round 23 of maximum 28.\\ \tt You have collectively accumulated 16 of 16 penalties.\\ \tt L'altro giocatore ha spostato un oggetto sulla sua griglia.\\ \tt \\ \tt \\ \tt Qual è il tuo prossimo comando? 
	  } 
	   } 
	   } 
	 & \\ 
 

    \theutterance \stepcounter{utterance}  

    & & & \multicolumn{4}{p{0.6\linewidth}}{\cellcolor[rgb]{1,0.85,0.72}{%
	\makecell[{{p{\linewidth}}}]{% 
	  \tt {\tiny [GM$\langle$B]}  
	 COMUNICA: Ho messo M@(7,6). Posso mettere O@(7,7) come prossimo passo; conferma il tuo stato attuale su x=7 dopo D e M così restiamo coordinati. 
	  } 
	   } 
	   } 
	  \\ 
 

    \theutterance \stepcounter{utterance}  

    & & & \multicolumn{2}{p{0.3\linewidth}}{\cellcolor[rgb]{0.95,0.95,0.95}{%
	\makecell[{{p{\linewidth}}}]{% 
	  \tt {\tiny [GM$|$GM]}  
	 COMUNICA: Ho messo M@(7,6). Posso mettere O@(7,7) come prossimo passo; conferma il tuo stato attuale su x=7 dopo D e M così restiamo coordinati. 
	  } 
	   } 
	   } 
	 & & \\ 
 

    \theutterance \stepcounter{utterance}  

    & & \multicolumn{4}{p{0.6\linewidth}}{\cellcolor[rgb]{0.9,0.9,0.9}{%
	\makecell[{{p{\linewidth}}}]{% 
	  \tt {\tiny [A$\langle$GM]}  
	 Spostato 'M' in (7,7) con successo. La tua griglia aggiornata appare così:\\ \tt ```\\ \tt  1234567\\ \tt ╔══╤════╗\\ \tt ║◌T│◌◌◌A║ 1\\ \tt ╟──┘◌◌┌─╢ 2\\ \tt ║S◌◌◌◌│◌║ 3\\ \tt ╟─┐◌┌─┼─╢ 4\\ \tt ║◌├─┤◌│O║ 5\\ \tt ╟─┤◌├─┤D║ 6\\ \tt ║◌│P│◌│M║ 7\\ \tt ╚═╧═╧═╧═╝\\ \tt ```\\ \tt You are currently playing round 23 of maximum 28.\\ \tt You have collectively accumulated 16 of 16 penalties.\\ \tt L'altro giocatore ha inviato questo messaggio:\\ \tt \\ \tt \\ \tt "Ho messo M@(7,6). Posso mettere O@(7,7) come prossimo passo; conferma il tuo stato attuale su x=7 dopo D e M così restiamo coordinati."\\ \tt \\ \tt \\ \tt Qual è il tuo prossimo comando? 
	  } 
	   } 
	   } 
	 & \\ 
 

    \theutterance \stepcounter{utterance}  

    & \multicolumn{4}{p{0.6\linewidth}}{\cellcolor[rgb]{0.8,1,0.9}{%
	\makecell[{{p{\linewidth}}}]{% 
	  \tt {\tiny [A$\rangle$GM]}  
	 COMUNICA: Stato mio su x=7: A@(7,1), D@(7,6), M@(7,7), O@(7,5). Le altre lettere non ancora su x=7. Puoi procedere con O@(7,7) sul tuo lato; da me è già occupato da M, quindi per O resterò su (7,5). Se vuoi, possiamo tentare P e S in altre colonne libere mantenendo l’ordine relativo sotto. 
	  } 
	   } 
	   } 
	 & & \\ 
 

    \theutterance \stepcounter{utterance}  

    & & & \multicolumn{2}{p{0.3\linewidth}}{\cellcolor[rgb]{0.95,0.95,0.95}{%
	\makecell[{{p{\linewidth}}}]{% 
	  \tt {\tiny [GM$|$GM]}  
	 COMUNICA: Stato mio su x=7: A@(7,1), D@(7,6), M@(7,7), O@(7,5). Le altre lettere non ancora su x=7. Puoi procedere con O@(7,7) sul tuo lato; da me è già occupato da M, quindi per O resterò su (7,5). Se vuoi, possiamo tentare P e S in altre colonne libere mantenendo l’ordine relativo sotto. 
	  } 
	   } 
	   } 
	 & & \\ 
 

    \theutterance \stepcounter{utterance}  

    & & \multicolumn{4}{p{0.6\linewidth}}{\cellcolor[rgb]{0.9,0.9,0.9}{%
	\makecell[{{p{\linewidth}}}]{% 
	  \tt {\tiny [GM$\rangle$B]}  
	 Il tuo messaggio è stato inoltrato all'altro giocatore.\\ \tt \\ \tt \\ \tt You are currently playing round 24 of maximum 28.\\ \tt You have collectively accumulated 16 of 16 penalties.\\ \tt L'altro giocatore ha inviato questo messaggio:\\ \tt \\ \tt \\ \tt "Stato mio su x=7: A@(7,1), D@(7,6), M@(7,7), O@(7,5). Le altre lettere non ancora su x=7. Puoi procedere con O@(7,7) sul tuo lato; da me è già occupato da M, quindi per O resterò su (7,5). Se vuoi, possiamo tentare P e S in altre colonne libere mantenendo l’ordine relativo sotto."\\ \tt \\ \tt \\ \tt Qual è il tuo prossimo comando? 
	  } 
	   } 
	   } 
	 & \\ 
 

    \theutterance \stepcounter{utterance}  

    & & & \multicolumn{4}{p{0.6\linewidth}}{\cellcolor[rgb]{1,0.85,0.72}{%
	\makecell[{{p{\linewidth}}}]{% 
	  \tt {\tiny [GM$\langle$B]}  
	 SPOSTA: O, (7,7) 
	  } 
	   } 
	   } 
	  \\ 
 

    \theutterance \stepcounter{utterance}  

    & & & \multicolumn{2}{p{0.3\linewidth}}{\cellcolor[rgb]{0.95,0.95,0.95}{%
	\makecell[{{p{\linewidth}}}]{% 
	  \tt {\tiny [GM$|$GM]}  
	 SPOSTA: O, (7,7) 
	  } 
	   } 
	   } 
	 & & \\ 
 

    \theutterance \stepcounter{utterance}  

    & & & \multicolumn{2}{p{0.3\linewidth}}{\cellcolor[rgb]{0.95,0.95,0.95}{%
	\makecell[{{p{\linewidth}}}]{% 
	  \tt {\tiny [GM$|$GM]}  
	 Spostato 'O' in (7,7) con successo. La tua griglia aggiornata appare così:\\ \tt ```\\ \tt  1234567\\ \tt ╔══╤════╗\\ \tt ║◌S│◌◌◌D║ 1\\ \tt ╟──┘◌◌┌─╢ 2\\ \tt ║◌◌◌◌◌│A║ 3\\ \tt ╟─┐◌┌─┼─╢ 4\\ \tt ║◌├─┤T│◌║ 5\\ \tt ╟─┤◌├─┤M║ 6\\ \tt ║P│◌│◌│O║ 7\\ \tt ╚═╧═╧═╧═╝\\ \tt ```\\ \tt  
	  } 
	   } 
	   } 
	 & & \\ 
 

    \theutterance \stepcounter{utterance}  

    & & \multicolumn{4}{p{0.6\linewidth}}{\cellcolor[rgb]{0.9,0.9,0.9}{%
	\makecell[{{p{\linewidth}}}]{% 
	  \tt {\tiny [A$\langle$GM]}  
	 Il tuo messaggio è stato inoltrato all'altro giocatore.\\ \tt \\ \tt \\ \tt You are currently playing round 24 of maximum 28.\\ \tt You have collectively accumulated 16 of 16 penalties.\\ \tt L'altro giocatore ha spostato un oggetto sulla sua griglia.\\ \tt \\ \tt \\ \tt Qual è il tuo prossimo comando? 
	  } 
	   } 
	   } 
	 & \\ 
 

    \theutterance \stepcounter{utterance}  

    & \multicolumn{4}{p{0.6\linewidth}}{\cellcolor[rgb]{0.8,1,0.9}{%
	\makecell[{{p{\linewidth}}}]{% 
	  \tt {\tiny [A$\rangle$GM]}  
	 COMUNICA: Aggiornamento finale mio: x=7 ha A@(7,1), O@(7,5), D@(7,6), M@(7,7). Se vuoi tentare ancora, posso provare a portare P sotto O in una colonna alternativa libera (per me x=6,y7 è libero? no, occupato da D a y6; x=5 è bloccato). Dimmi dove stai posizionando O,P,S,T sul tuo lato così decidiamo se fare un ultimo allineamento utile. 
	  } 
	   } 
	   } 
	 & & \\ 
 

    \theutterance \stepcounter{utterance}  

    & & & \multicolumn{2}{p{0.3\linewidth}}{\cellcolor[rgb]{0.95,0.95,0.95}{%
	\makecell[{{p{\linewidth}}}]{% 
	  \tt {\tiny [GM$|$GM]}  
	 COMUNICA: Aggiornamento finale mio: x=7 ha A@(7,1), O@(7,5), D@(7,6), M@(7,7). Se vuoi tentare ancora, posso provare a portare P sotto O in una colonna alternativa libera (per me x=6,y7 è libero? no, occupato da D a y6; x=5 è bloccato). Dimmi dove stai posizionando O,P,S,T sul tuo lato così decidiamo se fare un ultimo allineamento utile. 
	  } 
	   } 
	   } 
	 & & \\ 
 

    \theutterance \stepcounter{utterance}  

    & & \multicolumn{4}{p{0.6\linewidth}}{\cellcolor[rgb]{0.9,0.9,0.9}{%
	\makecell[{{p{\linewidth}}}]{% 
	  \tt {\tiny [GM$\rangle$B]}  
	 Spostato 'O' in (7,7) con successo. La tua griglia aggiornata appare così:\\ \tt ```\\ \tt  1234567\\ \tt ╔══╤════╗\\ \tt ║◌S│◌◌◌D║ 1\\ \tt ╟──┘◌◌┌─╢ 2\\ \tt ║◌◌◌◌◌│A║ 3\\ \tt ╟─┐◌┌─┼─╢ 4\\ \tt ║◌├─┤T│◌║ 5\\ \tt ╟─┤◌├─┤M║ 6\\ \tt ║P│◌│◌│O║ 7\\ \tt ╚═╧═╧═╧═╝\\ \tt ```\\ \tt You are currently playing round 25 of maximum 28.\\ \tt You have collectively accumulated 16 of 16 penalties.\\ \tt L'altro giocatore ha inviato questo messaggio:\\ \tt \\ \tt \\ \tt "Aggiornamento finale mio: x=7 ha A@(7,1), O@(7,5), D@(7,6), M@(7,7). Se vuoi tentare ancora, posso provare a portare P sotto O in una colonna alternativa libera (per me x=6,y7 è libero? no, occupato da D a y6; x=5 è bloccato). Dimmi dove stai posizionando O,P,S,T sul tuo lato così decidiamo se fare un ultimo allineamento utile."\\ \tt \\ \tt \\ \tt Qual è il tuo prossimo comando? 
	  } 
	   } 
	   } 
	 & \\ 
 

    \theutterance \stepcounter{utterance}  

    & & & \multicolumn{4}{p{0.6\linewidth}}{\cellcolor[rgb]{1,0.85,0.72}{%
	\makecell[{{p{\linewidth}}}]{% 
	  \tt {\tiny [GM$\langle$B]}  
	 COMUNICA: Sul mio lato x=7 ora: A@(7,3), D@(7,1), M@(7,6), O@(7,7). P è a (1,7), S a (2,1), T a (5,5). Posso tentare di portare P sotto O in colonna alternativa: candidato (6,7) libero sul mio lato. Se da te P può stare su x=6,y7, proviamo a mettere entrambi P@(6,7) come riferimento comune per P. Confermi? 
	  } 
	   } 
	   } 
	  \\ 
 

    \theutterance \stepcounter{utterance}  

    & & & \multicolumn{2}{p{0.3\linewidth}}{\cellcolor[rgb]{0.95,0.95,0.95}{%
	\makecell[{{p{\linewidth}}}]{% 
	  \tt {\tiny [GM$|$GM]}  
	 COMUNICA: Sul mio lato x=7 ora: A@(7,3), D@(7,1), M@(7,6), O@(7,7). P è a (1,7), S a (2,1), T a (5,5). Posso tentare di portare P sotto O in colonna alternativa: candidato (6,7) libero sul mio lato. Se da te P può stare su x=6,y7, proviamo a mettere entrambi P@(6,7) come riferimento comune per P. Confermi? 
	  } 
	   } 
	   } 
	 & & \\ 
 

    \theutterance \stepcounter{utterance}  

    & & \multicolumn{4}{p{0.6\linewidth}}{\cellcolor[rgb]{0.9,0.9,0.9}{%
	\makecell[{{p{\linewidth}}}]{% 
	  \tt {\tiny [A$\langle$GM]}  
	 Il tuo messaggio è stato inoltrato all'altro giocatore.\\ \tt \\ \tt \\ \tt You are currently playing round 25 of maximum 28.\\ \tt You have collectively accumulated 16 of 16 penalties.\\ \tt L'altro giocatore ha inviato questo messaggio:\\ \tt \\ \tt \\ \tt "Sul mio lato x=7 ora: A@(7,3), D@(7,1), M@(7,6), O@(7,7). P è a (1,7), S a (2,1), T a (5,5). Posso tentare di portare P sotto O in colonna alternativa: candidato (6,7) libero sul mio lato. Se da te P può stare su x=6,y7, proviamo a mettere entrambi P@(6,7) come riferimento comune per P. Confermi?"\\ \tt \\ \tt \\ \tt Qual è il tuo prossimo comando? 
	  } 
	   } 
	   } 
	 & \\ 
 

    \theutterance \stepcounter{utterance}  

    & \multicolumn{4}{p{0.6\linewidth}}{\cellcolor[rgb]{0.8,1,0.9}{%
	\makecell[{{p{\linewidth}}}]{% 
	  \tt {\tiny [A$\rangle$GM]}  
	 COMUNICA: Confermo P@(6,7) come riferimento comune. Sul mio lato (6,7) è libero, quindi dopo il tuo ok proverò a spostare P lì e ti confermerò l’esito. Poi possiamo valutare S e T in posizioni comuni se rimane tempo. 
	  } 
	   } 
	   } 
	 & & \\ 
 

    \theutterance \stepcounter{utterance}  

    & & & \multicolumn{2}{p{0.3\linewidth}}{\cellcolor[rgb]{0.95,0.95,0.95}{%
	\makecell[{{p{\linewidth}}}]{% 
	  \tt {\tiny [GM$|$GM]}  
	 COMUNICA: Confermo P@(6,7) come riferimento comune. Sul mio lato (6,7) è libero, quindi dopo il tuo ok proverò a spostare P lì e ti confermerò l’esito. Poi possiamo valutare S e T in posizioni comuni se rimane tempo. 
	  } 
	   } 
	   } 
	 & & \\ 
 

    \theutterance \stepcounter{utterance}  

    & & \multicolumn{4}{p{0.6\linewidth}}{\cellcolor[rgb]{0.9,0.9,0.9}{%
	\makecell[{{p{\linewidth}}}]{% 
	  \tt {\tiny [GM$\rangle$B]}  
	 Il tuo messaggio è stato inoltrato all'altro giocatore.\\ \tt \\ \tt \\ \tt You are currently playing round 26 of maximum 28.\\ \tt You have collectively accumulated 16 of 16 penalties.\\ \tt L'altro giocatore ha inviato questo messaggio:\\ \tt \\ \tt \\ \tt "Confermo P@(6,7) come riferimento comune. Sul mio lato (6,7) è libero, quindi dopo il tuo ok proverò a spostare P lì e ti confermerò l’esito. Poi possiamo valutare S e T in posizioni comuni se rimane tempo."\\ \tt \\ \tt \\ \tt Qual è il tuo prossimo comando? 
	  } 
	   } 
	   } 
	 & \\ 
 

    \theutterance \stepcounter{utterance}  

    & & & \multicolumn{4}{p{0.6\linewidth}}{\cellcolor[rgb]{1,0.85,0.72}{%
	\makecell[{{p{\linewidth}}}]{% 
	  \tt {\tiny [GM$\langle$B]}  
	 SPOSTA: P, (6,7) 
	  } 
	   } 
	   } 
	  \\ 
 

    \theutterance \stepcounter{utterance}  

    & & & \multicolumn{2}{p{0.3\linewidth}}{\cellcolor[rgb]{0.95,0.95,0.95}{%
	\makecell[{{p{\linewidth}}}]{% 
	  \tt {\tiny [GM$|$GM]}  
	 SPOSTA: P, (6,7) 
	  } 
	   } 
	   } 
	 & & \\ 
 

    \theutterance \stepcounter{utterance}  

    & & & \multicolumn{2}{p{0.3\linewidth}}{\cellcolor[rgb]{0.95,0.95,0.95}{%
	\makecell[{{p{\linewidth}}}]{% 
	  \tt {\tiny [GM$|$GM]}  
	 Penalità: (6,7) non è vuoto, ma contiene '│'.\\ \tt \\ \tt You have collectively accumulated 17 of 16 penalties.\\ \tt Riprova! 
	  } 
	   } 
	   } 
	 & & \\ 
 

    \theutterance \stepcounter{utterance}  

    & & & \multicolumn{2}{p{0.3\linewidth}}{\cellcolor[rgb]{0.95,0.95,0.95}{%
	\makecell[{{p{\linewidth}}}]{% 
	  \tt {\tiny [GM$|$GM]}  
	 Maximum number of penalties exceeded 
	  } 
	   } 
	   } 
	 & & \\ 
 

    \theutterance \stepcounter{utterance}  

    & & & \multicolumn{2}{p{0.3\linewidth}}{\cellcolor[rgb]{0.95,0.95,0.95}{%
	\makecell[{{p{\linewidth}}}]{% 
	  \tt {\tiny [GM$|$GM]}  
	 * success: False\\ \tt * lose: True\\ \tt * aborted: True\\ \tt -------\\ \tt * Shifts: 7.00\\ \tt * Max Shifts: 12.00\\ \tt * Min Shifts: 6.00\\ \tt * End Distance Sum: 19.24\\ \tt * Init Distance Sum: 19.71\\ \tt * Expected Distance Sum: 29.33\\ \tt * Penalties: 17.00\\ \tt * Max Penalties: 16.00\\ \tt * Rounds: 25.00\\ \tt * Max Rounds: 28.00\\ \tt * Object Count: 7.00\\ \tt  
	  } 
	   } 
	   } 
	 & & \\ 
 

    \theutterance \stepcounter{utterance}  

    & & & \multicolumn{2}{p{0.3\linewidth}}{\cellcolor[rgb]{0.95,0.95,0.95}{%
	\makecell[{{p{\linewidth}}}]{% 
	  \tt {\tiny [GM$|$GM]}  
	 * Main Score: 0.00\\ \tt \\ \tt -------\\ \tt * Distance Score: 0.18\\ \tt * Consistency Score: 0.86\\ \tt * Coverage Score: 0.31\\ \tt * Penalty Score: 0.00\\ \tt * Alternative Penalty Score: 0.00\\ \tt * Alternative Main Score: 0.00\\ \tt \\ \tt -------\\ \tt * Shifts: 7.00\\ \tt * Max Shifts: 12.00\\ \tt * Min Shifts: 6.00\\ \tt * End Distance Sum: 19.24\\ \tt * Init Distance Sum: 19.71\\ \tt * Expected Distance Sum: 29.33\\ \tt * Penalties: 17.00\\ \tt * Max Penalties: 16.00\\ \tt * Rounds: 25.00\\ \tt * Max Rounds: 28.00\\ \tt * Object Count: 7.00\\ \tt  
	  } 
	   } 
	   } 
	 & & \\ 
 

\end{supertabular}
}

\end{document}
